\chapter{Feynman diagram interpretation of bar-cobar duality}
\label{ch:v1-feynman}

\providecommand{\PhiFA}{\Phi^{\mathrm{FA}}}
\providecommand{\SpCh}{\mathrm{Sp}^{\mathrm{ch}}}

\begin{remark}[Bar is computational; Swiss-cheese is the explanatory
mechanism]
\label{rem:bar-computational-swiss-cheese-explanatory-feyn}
\index{bar-computational!corrected explanation|textbf}
\index{Swiss-cheese!explanatory mechanism|textbf}
\index{bar coalgebra!computational model|textbf}
Two distinctions govern this chapter and are load-bearing for every
later cross-volume reference.

\emph{Distinction 1 \textup{(}what is computed by what\textup{)}.}
The chiral bar coalgebra $B(A_b) = T^c(s^{-1}\bar A_b)$ is an
$E_1$-coassociative twisting coalgebra over the chiral associative
operad on the curve $X$ (Vol~I level~$2$ of the Beilinson tower).
It is a chain-level computational model that resolves $A_b$ and
presents the chiral Hochschild cochain complex
$\mathrm{C}^\bullet_{\mathrm{ch}}(A_b, A_b) \simeq
\mathrm{Hom}_{\mathrm{ch}}(B(A_b), A_b)$. The Feynman-diagram
interpretations developed below visualise this resolution.

\emph{Distinction 2 \textup{(}where the dimensional uplift comes
from\textup{)}.} The mechanism by which a $2$-real-dimensional chiral
algebra on $X$ produces a $3$-real-dimensional
holomorphic-topological bulk on $X \times \bR_{\ge 0}$ is the
\emph{Swiss-cheese pair} $(Z^{\mathrm{der}}_{\mathrm{ch}}(A_b),\, A_b)$
produced by the chiral Deligne--Tamarkin terminal-object construction:
the closed colour
$Z^{\mathrm{der}}_{\mathrm{ch}}(A_b) = \mathrm{C}^\bullet_{\mathrm{ch}}
(A_b, A_b)$ acts on the open colour $A_b$ through Voronov's
Swiss-cheese operations~\cite{Voronov99}; the open colour does not
act back on the closed colour. The extra topological dimension is the
universal normal-collar direction of a codimension-one boundary
condition, forced by the centre construction
\textup{(}Vol~I `chapters/theory/chiral\_center\_theorem.tex'
Theorem~\ref{thm:chiral-deligne-tamarkin}; cross-referenced by
Theorem~\ref{thm:swiss-cheese-chiral-DT-feynman} of the previous
chapter\textup{)}.

The bar direction by itself does not produce the dimensional uplift.
Drawings, BPHZ-shape recursions, and stable-graph formulas in this
chapter are computational presentations of the bar side; the
holographic mechanism passes through the Swiss-cheese pair, not
through the bar of the open algebra alone.  Cross-volume Vol~II
\texttt{thqg\_open\_closed\_realization.tex} carries the
Swiss-cheese realisation as theorem.
\end{remark}

\begin{remark}[Heuristic Feynman--bar mnemonic; \ClaimStatusHeuristic]
\label{rem:feynman-bar-bridge}
\index{Feynman diagram!bar-cobar bridge|textbf}
Let $\cA$ be a chiral algebra on a smooth curve~$X$ with OPE
$a(z)\,b(w) \sim \sum_n C^{ab}_n(w)\,(z-w)^{-n-1}$.
The ordered bar complex carries collision data on
Fulton--MacPherson compactifications.  A perturbative Feynman diagram
can be drawn on the same boundary stratification, but its Wick kernel is
not the Arnold logarithmic form.
\begin{center}
\renewcommand{\arraystretch}{1.25}
\begin{tabular}{lcl}
\textbf{Perturbative QFT} && \textbf{Bar complex $\barBch(\cA)$} \\
\hline
propagator $G(z_i,z_j) = (z_i-z_j)^{-1}$
 && logarithmic collision form
 $\eta_{ij} = d\log(z_i - z_j)$ \\
vertex factor $C^{ab}_n$
 && OPE residue $\Res_{z_i = z_j}[a(z_i)\,b(z_j)\,(z_i-z_j)^n]$ \\
UV divergence (short-distance singularity)
 && FM boundary stratum $D_S\subset\overline{C}_k(X)$, $|S|\ge 2$ \\
compactness/renormalization datum
 && logarithmic extension across $C_k(X) \hookrightarrow
 \overline{C}_k(X)$
\end{tabular}
\end{center}
This table is a mnemonic, not a chain equivalence.  Theorem~\ref{thm:bar-cobar-isomorphism-main}
and Theorem~\ref{thm:bar-nilpotency-complete} prove the algebraic
statement: the bar differential is the signed sum of chiral-product
residues along FM collision divisors, with the logarithmic forms
$\eta_{ij}$ supplying the bar-degree geometry and the Arnold relation
supplying the codimension-$2$ cancellation.  Proposition~\ref{prop:pole-decomposition}
separates the simple-pole bracket component from the scalar higher-pole
curvature packet: double-pole for Heisenberg and affine currents,
fourth-order before logarithmic shift for Virasoro. One does not
multiply an OPE pole by $\eta_{ij}$ and take a second ordinary residue.

A perturbative BV complex is compared with the bar complex only on the
proved loci stated below, and otherwise remains conjectural.  Graph loop
number is not the genus grading of the bar package:
Theorem~\ref{thm:loop-genus-correspondence} shows that ribbon-graph
genus also depends on face data.
\end{remark}

\index{Feynman diagram|textbf}

\section{Feynman diagrams and chiral residues}

\begin{definition}[Algebraic and perturbative diagram data]
\label{def:worldline-intro}
\index{propagator!Feynman}
\index{Feynman diagram!and bar complex}
The algebraic datum attached to a diagram on a curve~$X$ is a chiral
algebra~$\cA$, the logarithmic collision forms
$\eta_{ij}=d\log(z_i-z_j)$ at genus~$0$, and the OPE modes
$a_{(n)}b$ extracted by chiral residues.  On a fixed higher-genus curve
the local diagonal singularity is represented by the prime-form
logarithmic connection $d\log E(p_i,p_j)$ in a chosen trivialization of
the diagonal line; its residue along $p_i=p_j$ is still~$1$, while its
period corrections belong to the global curve and sewing data.  This
local statement defines the bar residue differential.  It does not by
itself choose a global Wick propagator.

A perturbative field-theory model supplies, in addition, Wick
contraction kernels and local vertices.  For a Feynman diagram~$\Gamma$
with $V$ vertices and $E$ internal edges, its formal configuration-space
amplitude has the shape
\[
A_\Gamma = \int_{C_V(X)}
\prod_{e=(i,j) \in E(\Gamma)} G(z_i,z_j)
\cdot \prod_{v \in V(\Gamma)} C_v
\cdot \prod_i dz_i,
\]
where $G$ is the chosen Wick kernel and $C_v$ are vertex factors.  The
bar form $\eta_{ij}$ records the collision divisor and its orientation;
it is not the Wick kernel unless a separate BV/bar comparison theorem
identifies the two roles in a particular free or anomaly-free model.
For the Cauchy kernel $G_C(z_i,z_j)=(z_i-z_j)^{-1}$ one has
$dG_C/G_C=-\eta_{ij}$.  For the Heisenberg current two-point kernel
$G_{\mathcal H}(z_i,z_j)=\kappa(z_i-z_j)^{-2}$ one has
$dG_{\mathcal H}/G_{\mathcal H}=-2\eta_{ij}$.  These are logarithmic
derivative identities, not identifications of kernels with collision
forms.
\end{definition}

\subsection{Tree vs.\ loop decomposition}

\begin{definition}[Loop number]
A Feynman diagram $\Gamma$ with $V$ vertices, $E$ edges, and $C$ connected
components has loop number:
\[
L(\Gamma) = E - V + C.
\]

This is the first Betti number $b_1(\Gamma)$ of the graph.
\end{definition}

\begin{remark}[Configuration space interpretation]
In a perturbative model, $L(\Gamma)$ counts independent graph cycles or
momentum integrations.  In the ordered bar complex, bar degree~$k$
counts $k+1$ field insertions before residue contraction.  These
gradings agree only after extra ribbon or worldline data have been
chosen; the FM boundary strata themselves are indexed by collision
patterns, not by physical loop order.
\end{remark}

\section{Off-shell/on-shell comparison}

The bar/cobar adjunction is algebraic.  A physical off-shell/on-shell
interpretation requires a BV/BRST model whose comparison map is part of
the data.

\begin{conjecture}[Feynman rules dictionary; \ClaimStatusConjectured]
\label{conj:physical-pairing}
\textup{(}\emph{quadrant} Open, \emph{presentation} bar / cobar pairing
\textup{(}computational level-$2$ presentation\textup{)},
\emph{level} 2 of the Beilinson tower with conjectural physical
identification, \emph{hypothesis package} BV/BRST factorisation model
whose local observables compare with $\cA$; identification of homotopies
with propagator kernels; this conjecture comparing pairings does not
itself produce the dimensional uplift, which is the
Theorem~\ref{thm:swiss-cheese-chiral-DT-feynman} mechanism.\textup{)}
For a BV/BRST factorization model whose local observables compare with
$\cA$, the bar-cobar construction is conjectured to encode the following
perturbative dictionary.  The bar complex carries logarithmic
collision forms on~$\overline{C}_n(X)$; the cobar complex carries the
Verdier-dual local-cohomology classes supported on collision strata; and
their pairing is the conjectural local model for the S-matrix pairing:
\[
\langle \omega_{\text{bar}}, K_{\text{cobar}} \rangle
= \int_{\overline{C}_n(X)}
\omega_{\text{bar}} \wedge \iota^* K_{\text{cobar}}
\;\sim\;
\langle \text{in} \,|\, S \,|\, \text{out} \rangle.
\]

\begin{center}
\begin{tabular}{lll}
\toprule
\textbf{Physical Object} & \textbf{Bar Complex} & \textbf{Cobar Complex} \\
\midrule
External leg & Boundary point & Marked point \\
Internal edge & Logarithmic form $\eta_{ij}$ & Local cohomology class on $D_{ij}$ \\
Vertex & Residue extraction & Comultiplication \\
Loop integration & Integration over $C_n$ & Trace over distributions \\
UV collision resolution & FM compactification & Support on boundary strata \\
IR datum & Global curve/moduli input & Dual global section \\
\bottomrule
\end{tabular}
\end{center}
\end{conjecture}

\begin{example}[Free boson propagator]
For the free boson $\alpha(z)$, the bar element
$\omega = \alpha(z_1) \otimes \alpha(z_2) \otimes \eta_{12}
\in \Omega^1(\overline{C}_2(X), \cA^{\boxtimes 2})$
carries the logarithmic collision form
$\eta_{12} = d\log(z_1-z_2)$.  The Verdier-dual
cobar element is the local cohomology class
$K \in H^1_{D_{12}}(\overline{C}_2(X), \cA^{!\boxtimes 2})$
supported on the diagonal $D_{12} = \{z_1 = z_2\}$. The pairing
$\langle \omega, K \rangle
= \operatorname{Res}_{D_{12}}(\eta_{12}) = 1$
is the logarithmic residue normalization.  The Wick contraction
$\kappa(z_1-z_2)^{-2}$ is a separate kernel; its scalar appears in the
bar differential through the curvature component
$\alpha_{(1)}\alpha=\kappa\mathbf{1}$.
\end{example}

\section{Higher operations from boundary strata}

\subsection{Virasoro residue packet}

\begin{example}[Virasoro ternary residue]\label{ex:virasoro-one-loop}
For the Virasoro algebra, the ternary bar operation is a genus-$0$
boundary operation on $\overline{M}_{0,4}$, not a genus-$1$ loop
integral.  A triangle drawing records the three collision channels.

\noindent\emph{Setup.}
\[T(z) = \sum_n L_n z^{-n-2}, \quad [L_m, L_n] = (m-n)L_{m+n} +
\frac{c}{12}m(m^2-1)\delta_{m+n,0}\]

OPE:
\[T(z)T(w) = \frac{c/2}{(z-w)^4} + \frac{2T(w)}{(z-w)^2} +
\frac{\partial T(w)}{z-w} + \text{reg}\]

\noindent\emph{Residue structure.}
The operation $m_3$ is represented, after a choice of homotopy transfer
data, by the signed sum of the three boundary residues of
$\overline{M}_{0,4} \cong \bP^1$:
\[m_3(T, T, T) = \sum_{\sigma \in \partial \overline{M}_{0,4}}
\epsilon_\sigma\,\mathrm{Res}_\sigma\bigl[\omega(T,T,T)\bigr].\]

Here $\omega$ is the logarithmic bar chain and
$\epsilon_\sigma$ is the FM boundary-orientation sign fixed by
Appendix~\ref{app:signs}.  In the channel $z_1\to z_2$, the local
coordinate $u=z_1-z_2$ gives the OPE modes
\[
T(z_2+u)T(z_2)
= \frac{c/2}{u^4}+\frac{2T(z_2)}{u^2}
 +\frac{\partial T(z_2)}{u}+{:}TT{:}(z_2)+\cdots .
\]
The fourth-order pole is the central Virasoro curvature mode
$T_{(3)}T=(c/2)\mathbf{1}$; the simple pole is the bracket mode
$T_{(0)}T=\partial T$.  The bar residue uses the full chiral product
and the logarithmic FM form.  It is not the ordinary residue of
$u^{-4}\,d\log u$.

The scalar genus-$1$ invariant is read instead from the modular MC
element:
\[
F_1(\mathrm{Vir}_c)=\frac{\kappa(\mathrm{Vir}_c)}{24}
=\frac{c}{48}.
\]
Thus the same central charge controls both the genus-$0$ pole
decomposition and the genus-$1$ scalar obstruction, but the ternary
residue is not itself a one-loop theorem.
\end{example}

\subsection{Higher loops and factorization}

\begin{theorem}[\texorpdfstring{$A_\infty$}{A-infinity} constraint formula; \ClaimStatusProvedHere]
\label{thm:ainfty-constraint-formula}
\textup{(}\emph{quadrant} Open, \emph{presentation} chart $A_b$,
\emph{level} 1 of the Beilinson tower, \emph{hypothesis package}
uncurved $\Ainf$-structure with $m_0 = 0$ in Keller signs.\textup{)}
For an uncurved $A_\infty$ algebra written in Keller signs, the
operations satisfy the
constraint
\[
m_1 \circ m_n =
-\!\!\sum_{\substack{r+s+t=n,\ s\geq 1\\ (r,s,t) \neq (0,n,0)}}
(-1)^{rs+t}
m_{r+1+t}(\mathrm{id}^{\otimes r}\otimes m_s\otimes
\mathrm{id}^{\otimes t}).
\]

This constrains $m_1(m_n)$ (not $m_n$ itself) in terms of lower operations,
and gives the algebraic Ward identity to which BPHZ recursion can be
compared.
\end{theorem}

\begin{proof}
This is a rewriting of the uncurved $A_\infty$ relations
(Keller convention):
\[
\sum_{\substack{r+s+t=n\\ s\geq 1}} (-1)^{rs+t}
m_{r+1+t}(\mathrm{id}^{\otimes r} \otimes m_s \otimes
\mathrm{id}^{\otimes t}) = 0.
\]

Isolating the $r = t = 0$, $s = n$ term, which is $m_1(m_n)$, gives the result.
To define $m_n$ recursively by homotopy transfer one also needs a
contracting homotopy $h$ with
$h m_1 + m_1 h = \mathrm{id} - \iota\pi$; the displayed identity is a
consistency relation among already-defined operations.
\end{proof}

\begin{remark}[Curvature and the MC half-factor]
\label{rem:feynman-curvature-qme-half}
In the curved case the $s=0$ terms are present and contain the
curvature operation $m_0$.  On the chiral bar side,
Proposition~\ref{prop:pole-decomposition} identifies the corresponding
curvature component with scalar higher-pole OPE modes, double-pole in
the affine and Heisenberg examples and fourth-order in the Virasoro
OPE before the logarithmic collision shift.  The modular
convolution algebra packages the same data as a Maurer--Cartan element
\[
D_\cA\Theta_\cA+\tfrac{1}{2}[\Theta_\cA,\Theta_\cA]=0.
\]
The factor $\tfrac{1}{2}$ is the standard dg-Lie normalization of the
quadratic bracket term, matching the BV quantum master equation
normalization recorded in Appendix~\ref{app:signs}.  It is not a graph
symmetry factor chosen separately for each diagram.
\end{remark}

\section{Graph complexes and Kontsevich formality}

\subsection{The graph complex}

\begin{definition}[Decorated graph complex convention]
The graph complex used in this section is the graded vector space
$\mathrm{GC}_n(\cA)$ generated by finite graphs with $n$ labelled
external legs, unlabeled internal vertices, and vertex decorations in
the local operations of~$\cA$.  Its differential is the signed sum of
the edge contractions and vertex splittings allowed by the relevant
operad; in Kontsevich's original admissible-graph model this same
boundary operation is usually written in vertex-splitting form.
The grading retained here is:
\begin{itemize}
\item external arity $n$;
\item graph loop number $L(\Gamma)=b_1(\Gamma)$;
\item total stable genus after a ribbon or modular decoration is
 supplied.
\end{itemize}
\end{definition}

\begin{conjecture}[Bar complex and graph complex; \ClaimStatusConjectured]
\label{conj:v1-bar-graph-complex}
\textup{(}\emph{quadrant} Open, \emph{presentation} bar / twisting
coalgebra in the chain-level Kontsevich graph-complex model,
\emph{level} 2 of the Beilinson tower, \emph{hypothesis package}
Kontsevich admissible-graph signs in the chiral specialisation; FM
collision strata; OPE decorations; logarithmic forms $\eta_{ij}$;
the conjecture is at the level of the bar computational model and
does not displace the chiral Deligne--Tamarkin
\textup{(}Theorem~\ref{thm:swiss-cheese-chiral-DT-feynman}\textup{)}
explanatory mechanism for the dimensional uplift.\textup{)}
The chiral bar complex admits a quasi-isomorphism to the
$\cA$-decorated Kontsevich graph complex:
$\barB^{\mathrm{ch}}(\cA) \simeq \mathrm{GC}(\cA)$,
where vertices correspond to field insertions, edges to logarithmic
forms~$\eta_{ij}$, and the graph differential matches the signed FM
residue differential after the OPE decorations are inserted.
\end{conjecture}

\subsection{Kontsevich's formality and chiral algebras}

\begin{theorem}[Kontsevich formality; \ClaimStatusProvedElsewhere]
\label{thm:kontsevich-formality-feynman}
\textup{(}\emph{quadrant} smooth manifold $M$ \textup{(}not chiral; cited
for analogue\textup{)}, \emph{presentation} polydifferential Hochschild
cochains, \emph{level} bulk Hochschild,
\emph{hypothesis package} $C^\infty$ smoothness; admissible-graph
configuration-space convergence by~\cite{Kon03}.\textup{)}
For a smooth manifold $M$, the dg Lie algebra
$D_{\mathrm{poly}}(M)$ of polydifferential Hochschild cochains of
$C^\infty(M)$ is $L_\infty$ quasi-isomorphic to the graded Lie algebra
$T_{\mathrm{poly}}(M)$ of polyvector fields with the Schouten bracket:
\[
\mathcal U_{\mathrm{Kont}}\colon
T_{\mathrm{poly}}(M)\xrightarrow{\simeq}D_{\mathrm{poly}}(M).
\]
Equivalently, the Hochschild deformation complex is formal with
cohomology $T_{\mathrm{poly}}(M)$, and the Taylor coefficients of
$\mathcal U_{\mathrm{Kont}}$ are configuration-space integrals over
admissible graphs~\cite{Kon03}. The chiral analogue at the level of
Swiss-cheese operations is recorded in
Theorem~\ref{thm:swiss-cheese-chiral-DT-feynman}; the present formality
theorem cites the $C^\infty$ topological situation, which is the model
to which the chiral Deligne--Tamarkin theorem is the chiral upgrade.
\end{theorem}

\begin{remark}[Chiral analogue of Kontsevich formality]
\label{rem:kontsevich-graphs}
\label{sec:connections-other-feynman}
The bar-cobar construction on a curve~$X$ is analogous to
Kontsevich formality~\cite{Kon99}: angle forms are replaced by the
Arnold logarithmic forms, and graph boundary strata are replaced by FM
collision strata.  The actual theorem in this monograph is the
bar-cobar adjunction and its residue differential; a Kontsevich-style
configuration-integral formality morphism for arbitrary chiral
algebras is not proved here.
\end{remark}

\begin{conjecture}[Disk-local perturbative/FM comparison;
 \ClaimStatusConjectured]\label{conj:v1-disk-local-perturbative-fm}
\textup{(}\emph{quadrant} Open, \emph{presentation} chart $A_b$ on a
holomorphic disk $U$, \emph{level} 1 of the Beilinson tower,
\emph{hypothesis package} disk chart $U \subset X$;
holomorphic-topological perturbative field theory whose local chiral
algebra model is $\cA|_U$; Gaiotto--Kulp--Wu~\cite{GKW24} regularised
BRST-anomaly brackets; FM signs / orientations of
Appendix~\ref{app:signs}.\textup{)}
Let $U \subset X$ be a holomorphic disk chart on a curve, and suppose a
holomorphic-topological perturbative field theory on~$U$ produces the
same local chiral algebra model $\cA|_U$.\linebreak[1]
Then the regularized perturbative BRST-anomaly brackets
of Gaiotto--Kulp--Wu~\cite{GKW24} agree, on that disk chart, with
the local bar operations cut out by the Fulton--MacPherson collision
strata:
\begin{enumerate}[label=\textup{(\roman*)}]
\item the binary perturbative bracket agrees with the bar residue/OPE
 operation on~$C_2(U)$;
\item the ternary perturbative bracket agrees with the sum of bar
 residues on the three codimension-$1$ collision strata of
 $\overline{C}_3(U)$, hence with the local $m_3$-operation; and
\item the quadratic identities satisfied by those perturbative brackets
 agree with the Stokes relations among the corresponding
 codimension-$1$ FM strata after passage to the compactification.
\end{enumerate}
\end{conjecture}

\begin{remark}[Disk-local comparison packet]\label{rem:mc5-local-packet}
Conjecture~\ref{conj:v1-disk-local-perturbative-fm} isolates the first
finite comparison locus: a binary collision calculation on~$C_2$, a
ternary calculation on the three boundary divisors of
$\overline{C}_3$, and the identification of the perturbative
quadratic/Wess--Zumino relation with the FM Stokes relation.  The
algebraic Stokes relation is proved inside the bar complex; its
identification with the perturbative relation is the conjectural
content.
\end{remark}

\begin{proposition}[Binary-ternary reduction for the disk-local packet;
 \ClaimStatusConditional]\label{prop:disk-local-binary-ternary-reduction}
Fix a holomorphic disk chart $U \subset X$, and interpret the
perturbative BRST-anomaly brackets and the local bar operations on the
same local chain model underlying~$\cA|_U$.
Assume:
\begin{enumerate}[label=\textup{(\roman*)}]
\item the binary perturbative bracket agrees with the local bar
 residue/OPE operation on~$C_2(U)$; and
\item the ternary perturbative bracket agrees with the local
 $m_3$-operation, equivalently with the sum of the three
 codimension-$1$ residues on~$\overline{C}_3(U)$.
\end{enumerate}
Then the quadratic/Wess--Zumino identities on the perturbative side
agree with the quadratic $A_\infty$/Stokes relation on the bar side,
provided the perturbative orientations are identified with the FM
boundary orientations of Appendix~\ref{app:signs}.  Under this sign
identification there is no third independent local coefficient
calculation beyond the binary and ternary comparisons.
\end{proposition}

\begin{proof}
On the perturbative side, Gaiotto--Kulp--Wu~\cite{GKW24} prove that
the regularized BRST-anomaly brackets satisfy quadratic
Wess--Zumino consistency relations.

On the bar side, the corresponding degree-$3$ consistency relation is
already built into the local $A_\infty$ structure: by
Theorem~\ref{thm:ainfty-constraint-formula} specialized to $n=3$, the
bar operations satisfy the quadratic relation expressing $m_1m_3$ in
terms of composites of lower operations.  Geometrically this is the
Stokes relation among the codimension-$1$ boundary strata of the
compactified three-point configuration space, with signs fixed by the
FM orientation system.

After the binary and ternary operations are identified on the common
local chain model and the orientation systems are matched, both sides
express the same quadratic relation.  The comparison therefore has no
third independent local coefficient.
\end{proof}

\begin{remark}[Application to the disk-local FM comparison]
Applied to Conjecture~\ref{conj:v1-disk-local-perturbative-fm},
Proposition~\ref{prop:disk-local-binary-ternary-reduction} shows that
the only independent local input is the binary comparison on $C_2(U)$
and the ternary comparison on $\overline{C}_3(U)$.
\end{remark}

\begin{corollary}[Anomaly-free genus-\texorpdfstring{$0$}{0} collapse of the local packet;
 \ClaimStatusConditional]\label{cor:disk-local-ternary-on-brstbar-locus}
On the anomaly-free genus-$0$ BRST/bar locus of
Theorem~\ref{thm:brst-bar-genus0}, the binary comparison in
Proposition~\ref{prop:disk-local-binary-ternary-reduction} is already
proved. Equivalently, on that theorem surface the remaining
independent local MC5 check is the ternary comparison on
$\overline{C}_3$.
\end{corollary}

\begin{proof}
Theorem~\ref{thm:brst-bar-genus0} identifies the genus-$0$ BRST complex
with the genus-$0$ bar complex by a chain map
\(\Phi\) to \(\barB^{\mathrm{ch}}_0(\cA_{\mathrm{tot}})\).
Remark~\ref{rem:brst-bar-explicit} makes the low-degree content of that
map concrete: physical vertex operators map to bar generators, and the
ghost bilinear maps to the logarithmic collision form
\(\eta_{12}=d\log(z_1-z_2)\).
Thus the bar-generator and logarithmic-collision data underlying the
binary local packet are already identified on this proved genus-$0$
locus.

By Proposition~\ref{prop:disk-local-binary-ternary-reduction}, once the
binary packet is in hand the only remaining independent local
comparison is the ternary one.
\end{proof}

\begin{proposition}[Two-channel reduction after compactifying the
 ternary packet; \ClaimStatusProvedHere]
\label{prop:compactified-ternary-two-channel}
\textup{(}\emph{quadrant} Open, \emph{presentation} chart $A_b$,
\emph{level} 1 of the Beilinson tower,
\emph{hypothesis package} disk chart $U \subset X$; perturbative and
bar ternary forms representable by logarithmic $1$-forms on
$\overline{\mathcal{M}}_{0,4} \cong \bP^1$ with poles only at the
three boundary points; residue theorem on
$\bP^1$.\textup{)}
Fix a holomorphic disk chart $U \subset X$, and suppose the
perturbative ternary bracket and the local bar $m_3$-operation are
represented by logarithmic $1$-forms on
$\overline{M}_{0,4} \cong \bP^1$ whose only poles lie at the three
boundary points corresponding to the channels
$(12|3)$, $(23|1)$, and $(13|2)$.
Then equality of the perturbative and bar residues in any two of those
channels
implies comparison in the third channel as well.
Equivalently, after compactification the ternary packet has only two
independent boundary-channel checks.
\end{proposition}

\begin{proof}
Let $\omega_{\mathrm{pert}}$ and $\omega_{\mathrm{bar}}$ denote the two
logarithmic $1$-forms on $\overline{M}_{0,4} \cong \bP^1$, and set
\[
\delta = \omega_{\mathrm{pert}} - \omega_{\mathrm{bar}}
\in H^0\!\left(\bP^1,\,
\Omega^1_{\bP^1}\bigl(\log\{0,1,\infty\}\bigr)\right).
\]
If the perturbative and bar residues agree in two channels, then two of
the three residues of $\delta$ vanish. By the residue theorem on~$\bP^1$,
the sum of all three residues is zero, so the third residue vanishes as
well. Therefore $\delta$ has zero residue at every boundary point, so
it is a holomorphic $1$-form on~$\bP^1$. But
$H^0(\bP^1,\Omega^1_{\bP^1}) = 0$, hence $\delta = 0$.
In particular, the third boundary residue is forced by the first two.
\end{proof}

\begin{remark}[Application to the disk-local ternary packet]
Proposition~\ref{prop:compactified-ternary-two-channel} applies in
particular to the compactified ternary comparison sought in
Conjecture~\ref{conj:v1-disk-local-perturbative-fm}. Once the local
ternary forms are constructed on $\overline{M}_{0,4}$, only two
boundary-channel checks remain independent.
\end{remark}

\begin{corollary}[Genus-\texorpdfstring{$0$}{0} post-compactification ternary target;
 \ClaimStatusConditional]
\label{cor:genus0-compactified-ternary-two-channel}
On the anomaly-free genus-$0$ BRST/bar locus of
Theorem~\ref{thm:brst-bar-genus0}, once the compactification/Stokes
upgrade has been achieved, the remaining local MC5 comparison reduces
to two ternary boundary-channel checks on
$\overline{M}_{0,4} \cong \bP^1$.
\end{corollary}

\begin{proof}
Corollary~\ref{cor:disk-local-ternary-on-brstbar-locus} reduces the
local MC5 problem on this proved genus-$0$ locus to the ternary comparison
alone. Proposition~\ref{prop:compactified-ternary-two-channel} then
reduces that compactified ternary comparison to any two of the three
boundary channels.
\end{proof}

\begin{corollary}[Standard-chart form of the remaining genus-\texorpdfstring{$0$}{0}
 packet; \ClaimStatusConditional]
\label{cor:genus0-standard-chart-two-residues}
Choose the standard coordinate on
$\overline{M}_{0,4} \cong \bP^1$ with boundary points
$0,1,\infty$. On the anomaly-free genus-$0$ BRST/bar locus of
Theorem~\ref{thm:brst-bar-genus0}, after the compactification/Stokes
upgrade the remaining local MC5 comparison is exactly the equality of
the compactified ternary residues at $0$ and~$1$; the residue at
$\infty$ is then forced.
\end{corollary}

\begin{proof}
By Corollary~\ref{cor:genus0-compactified-ternary-two-channel}, any two
boundary-channel checks suffice. On the standard
$0,1,\infty$ compactification of $\overline{M}_{0,4}$, choose
the channels at $0$ and~$1$ as the independent pair. Then
Proposition~\ref{prop:compactified-ternary-two-channel} forces the
residue at~$\infty$.
\end{proof}

\begin{proposition}[Standard logarithmic basis on
 \texorpdfstring{$\overline{M}_{0,4}$}{M_0,4}; \ClaimStatusProvedHere]
\label{prop:m04-standard-log-basis}
\textup{(}\emph{quadrant} purely combinatorial
\textup{(}independent of chiral algebra\textup{)},
\emph{presentation} logarithmic forms on $\bP^1$,
\emph{level} 5 \textup{(}genus-$0$ stratum on
$\overline{\mathcal{M}}_{0,4}$\textup{)},
\emph{hypothesis package} simple poles only at $\{0, 1, \infty\}$ on
$\overline{\mathcal{M}}_{0,4} \cong \bP^1$;
$H^0(\bP^1, \Omega^1_{\bP^1}) = 0$.\textup{)}
Let $t$ be the standard coordinate on
$\overline{M}_{0,4} \cong \bP^1$ with boundary points
$0,1,\infty$. Then every logarithmic $1$-form on
$\bP^1$ with poles only at $\{0,1,\infty\}$ admits a unique
expression
\[
\omega = A\, d\log t + B\, d\log(1-t)
\]
for scalars $A,B \in \bC$.
Equivalently, the residues at $0$ and $1$ determine the form
uniquely, and the residue at $\infty$ equals $-A-B$.
\end{proposition}

\begin{proof}
Write $\omega = f(t)\,dt$ with $f(t)$ a rational function on~$\bP^1$
having at most simple poles at $0$, $1$, and $\infty$.
Thus
\[
f(t)=\frac{a}{t}+\frac{b}{t-1}+c
\]
for some scalars $a,b,c$. The condition that $\omega$ have at most a
simple pole at $\infty$ forces $c=0$: in the coordinate $u=1/t$ one
has
\[
dt = -u^{-2}\,du,\qquad
\left(\frac{a}{t}+\frac{b}{t-1}+c\right)dt
= \left(\frac{a}{u^{-1}}+\frac{b}{u^{-1}-1}+c\right)(-u^{-2}\,du),
\]
and a nonzero constant term $c$ would contribute a double pole at
$u=0$.
Hence
\[
\omega = a\,\frac{dt}{t}+b\,\frac{dt}{t-1}
= a\, d\log t - b\, d\log(1-t).
\]
Setting $A=a$ and $B=-b$ gives existence.

For uniqueness, if
$A\, d\log t + B\, d\log(1-t)=0$, then the residues at
$t=0$ and $t=1$ are respectively $A$ and $B$, so $A=B=0$.
The residue at $\infty$ is minus the sum of the finite residues, hence
$-A-B$.
\end{proof}

\begin{corollary}[Two-coefficient form of the remaining genus-\texorpdfstring{$0$}{0}
 compactified packet; \ClaimStatusConditional]
\label{cor:genus0-two-coefficient-packet}
On the anomaly-free genus-$0$ BRST/bar locus of
Theorem~\ref{thm:brst-bar-genus0}, after the compactification/Stokes
upgrade the remaining local MC5 comparison is exactly the vanishing of
two scalar coefficients. More precisely, if
\[
\delta = \omega_{\mathrm{pert}} - \omega_{\mathrm{bar}}
\in H^0\!\left(\bP^1,\,
\Omega^1_{\bP^1}\bigl(\log\{0,1,\infty\}\bigr)\right),
\]
then uniquely
\[
\delta = A\, d\log t + B\, d\log(1-t),
\]
and the compactified genus-$0$ comparison is equivalent to
$A=B=0$.
\end{corollary}

\begin{proof}
Corollary~\ref{cor:genus0-standard-chart-two-residues} identifies the
remaining task with equality of the residues at $0$ and $1$ in the
standard chart. Proposition~\ref{prop:m04-standard-log-basis}
identifies those two residues with the coefficients of
$d\log t$ and $d\log(1-t)$ respectively.
\end{proof}

\begin{definition}[Standard genus-\texorpdfstring{$0$}{0} coefficient packet]
\label{def:genus0-mc5-coefficient-packet}
For
\[
\omega \in H^0\!\left(\bP^1,\,
\Omega^1_{\bP^1}\bigl(\log\{0,1,\infty\}\bigr)\right),
\]
write uniquely
\[
\omega=\mathsf{a}(\omega)\,d\log t+\mathsf{b}(\omega)\,d\log(1-t)
\]
as in Proposition~\ref{prop:m04-standard-log-basis}. For the
compactified ternary perturbative and bar forms, write
\[
\omega_{\mathrm{pert}}
= \mathsf{a}_{\mathrm{pert}}\,d\log t
 + \mathsf{b}_{\mathrm{pert}}\,d\log(1-t),
\qquad
\omega_{\mathrm{bar}}
= \mathsf{a}_{\mathrm{bar}}\,d\log t
 + \mathsf{b}_{\mathrm{bar}}\,d\log(1-t).
\]
Call the pair of equalities
\[
\mathsf{a}_{\mathrm{pert}}=\mathsf{a}_{\mathrm{bar}},
\qquad
\mathsf{b}_{\mathrm{pert}}=\mathsf{b}_{\mathrm{bar}}
\]
the \emph{standard genus-$0$ MC5 coefficient packet}.
\end{definition}

\begin{corollary}[Named coefficient form of the remaining genus-\texorpdfstring{$0$}{0}
 compactified packet; \ClaimStatusConditional]
\label{cor:genus0-named-coefficient-packet}
On the anomaly-free genus-$0$ BRST/bar locus of
Theorem~\ref{thm:brst-bar-genus0}, after the compactification/Stokes
upgrade the remaining local MC5 comparison is exactly the standard
genus-$0$ MC5 coefficient packet of
Definition~\ref{def:genus0-mc5-coefficient-packet}.
\end{corollary}

\begin{proof}
Corollary~\ref{cor:genus0-two-coefficient-packet} identifies the
compactified genus-$0$ comparison with the vanishing of the two
coefficients of
\(\delta=\omega_{\mathrm{pert}}-\omega_{\mathrm{bar}}\).
This is equivalent to the two equalities
\[
\mathsf{a}_{\mathrm{pert}}=\mathsf{a}_{\mathrm{bar}},
\qquad
\mathsf{b}_{\mathrm{pert}}=\mathsf{b}_{\mathrm{bar}}.
\]
\end{proof}

\section{\texorpdfstring{The $m_k$ operations and diagram heuristics}{The m k operations and diagram heuristics}}
\label{sec:mk-feynman-complete}

\subsection{\texorpdfstring{The algebraic $m_k$ family}{The algebraic m k family}}

\begin{definition}[Diagrammatic reading of the \texorpdfstring{$m_k$}{m_k} family]
\label{def:mk-family-feynman}
The $A_\infty$ operations
$m_k\colon \mathcal A^{\otimes k}\to\mathcal A$ extend to
coderivations on the bar complex.  Under a contraction
$(H,\iota,\pi,h)$, the transferred operation $m_k^H$ is a sum over
planar rooted trees with $k$ leaves.  The tree edges are homotopy
insertions $h$; they are not physical loop edges unless a perturbative
model supplies such an identification.

\begin{center}
\begin{tabular}{|>{\raggedright\arraybackslash}p{1.2cm}|>{\raggedright\arraybackslash}p{4.4cm}|>{\raggedright\arraybackslash}p{5.7cm}|}
\hline
\textbf{$k$} & \textbf{Algebraic role} & \textbf{Permitted diagram reading} \\
\hline
$m_0$ & Curvature term & Vacuum or tadpole insertion in a curved perturbative model \\
\hline
$m_1$ & Differential & BRST differential when a BV/BRST model is chosen \\
\hline
$m_2$ & Binary chiral product/OPE residue & Tree vertex \\
\hline
$m_3$ & Ternary associator from $\partial\overline{M}_{0,4}$ & Triangle-shaped contact term; not a genus-$1$ loop by itself \\
\hline
$m_4$ & Pentagon transfer sum & Five tree terms of the Stasheff associahedron \\
\hline
$m_k$ & Sum over $k$-leaf transfer trees & A Feynman amplitude only after an external BV/bar comparison \\
\hline
\end{tabular}
\end{center}
\end{definition}

\begin{theorem}[Graph loop number and ribbon genus; \ClaimStatusProvedElsewhere{} \cite{costello-renormalization}]
\label{thm:loop-genus-formula}
\textup{(}\emph{quadrant} graph-combinatorial \textup{(}cited\textup{)},
\emph{presentation} ribbon-graph CW datum, \emph{level} 5 \textup{(}genus
filtration on $\overline{\mathcal{M}}_{g,n}$\textup{)},
\emph{hypothesis package} connected ribbon graph with $V$ vertices,
$E$ edges, $F$ faces; first Betti number; CW decomposition of the
boundary-capped surface.\textup{)}
For a connected graph $\Gamma$ with $V$ vertices and $E$ internal edges,
the graph loop number is
\[
\ell(\Gamma) = E - V + 1 = b_1(\Gamma).
\]
where $b_1$ is the first Betti number of $\Gamma$ viewed as a 1-complex.

When $\Gamma$ is equipped with a ribbon structure, it thickens to a
surface with $F$ faces, and
\[
2-2g=V-E+F,\qquad
g=\frac{\ell+1-F}{2}.
\]
Thus loop number alone does not determine the genus.

For chiral algebras, the genus-$0$ transferred operations are indexed
by stable trees in $\partial\overline{M}_{0,k+1}$.  The number
$k-2$ is the number of internal homotopy edges in a binary transfer
tree, not the genus of a Riemann surface.
\end{theorem}

\begin{proof}
For a connected graph $\Gamma$ with $V$ vertices and $E$ edges,
the first Betti number is $b_1(\Gamma) = E - V + 1$. This equals the
number of independent graph cycles.

For a ribbon graph $\Gamma$ thickened to a surface $\Sigma$ with $F$ boundary components, cap the boundary to obtain a closed surface $\bar{\Sigma}$ of genus $g$. The CW structure on $\bar{\Sigma}$ has $V$ vertices, $E$ edges, and $F$ faces (the capping disks), so:
\[\chi(\bar{\Sigma}) = V - E + F = 2 - 2g\]
Hence $g=(\ell+1-F)/2$.  The boundary of
$\overline{M}_{0,k+1}$ is stratified by stable trees, and each
codimension-one face records a composition of lower operations.  Stokes'
theorem on these faces gives the $A_\infty$ relations.
\end{proof}

\subsection{\texorpdfstring{$m_2$: binary collision}{m 2: binary collision}}

\begin{example}[Binary product = classical OPE]
\label{ex:m2-tree-level}
On a degree-$1$ bar chain
$\phi_1(z_1)\otimes\phi_2(z_2)\otimes\eta_{12}$, the binary
collision operation is the Poincar\'e residue of the chiral product:
\[
d_{\mathrm{res}}\bigl(\phi_1\otimes\phi_2\otimes\eta_{12}\bigr)
=\operatorname{Res}_{D_{12}}\bigl[\mu(\phi_1,\phi_2)\otimes\eta_{12}\bigr]
\]
The Poincar\'e residue sends the logarithmic factor
$\eta_{12}=d\log(z_1-z_2)$ to~$1$; the chiral product supplies the OPE
modes as a collision packet, with pole order recorded by the local
normal coordinate and by the chiral product.  On the simple-pole
bracket lane this reduces to
$m_2^{\mathrm{br}}(\phi_1,\phi_2)=(\phi_1)_{(0)}\phi_2$.
A physical three-point scattering amplitude is an additional
realization, not the definition of $m_2$.
\end{example}

\begin{remark}[OPE interpretation]\label{rem:witten-m2}
Write
\[
\phi_1(z)\phi_2(w)\sim
\sum_{n\geq 0}\frac{(\phi_1)_{(n)}\phi_2(w)}{(z-w)^{n+1}}.
\]
The bar residue extracts this mode packet through the chiral product.
The coefficient of the simple-pole projection is the bracket
structure constant; higher-pole modes enter the curvature and
secondary components of Proposition~\ref{prop:pole-decomposition},
with scalar higher-pole modes giving the scalar curvature terms.
Path-integral language may realize the same matrix coefficients as
tree-level vertices after a BV/bar comparison has been chosen.
\end{remark}

\subsection{\texorpdfstring{$m_3$: ternary boundary residue}{m 3: ternary boundary residue}}

\begin{example}[Ternary operation = three boundary channels]
\label{ex:m3-one-loop}
After choosing homotopy-transfer data, the genus-$0$ ternary
operation has the boundary form
\[
m_3(\phi_1,\phi_2,\phi_3)
=\sum_{\sigma\in\{12|3,\,23|1,\,13|2\}}
\epsilon_\sigma\,
\operatorname{Res}_{D_\sigma}
\bigl[\mu_\sigma(\phi_1,\phi_2,\phi_3)\otimes\omega_\sigma\bigr],
\]
where $D_\sigma\subset\partial\overline{M}_{0,4}$ is the boundary
point for the channel, $\omega_\sigma$ is the corresponding
logarithmic FM form, and $\mu_\sigma$ is the iterated chiral product
in that channel.

The boundary has the three channels $(12|3)$, $(23|1)$, and
$(13|2)$.  A triangle diagram is a drawing of these three channels; it
does not introduce a genus-$1$ integration variable in the bar complex.
\end{example}

\begin{computation}[Virasoro \texorpdfstring{$m_3$}{m_3} residue modes]
\label{comp:v1-virasoro-m3}
For the Virasoro algebra with stress tensor $T(z)$:

\emph{OPE:}
\[T(z)T(w) = \frac{c/2}{(z-w)^4} + \frac{2T(w)}{(z-w)^2} + \frac{\partial T(w)}{z-w}
+ \text{reg}\]

\emph{Boundary channels.}
The boundary of $\overline{M}_{0,4}$ has three components
corresponding to the three collision orders:

\[
m_3(T,T,T)=
\sum_{\sigma\in\{12|3,\,23|1,\,13|2\}}
\epsilon_\sigma\,\operatorname{Res}_{D_\sigma}
\bigl[\mu_\sigma(T,T,T)\otimes\omega_\sigma\bigr].
\]

\emph{First term.}
Collide $z_1 \to z_2$ first. Setting $u = z_1 - z_2$, the logarithmic
bar form near $D_{12}$ has the factor $d\log u$:
\[
\mu(T(z_2+u),T(z_2))\otimes T(z_3)\otimes d\log u\wedge\omega_{23}.
\]
The OPE gives $T(z_2+u)T(z_2) = c/2 \cdot u^{-4} + 2T(z_2) \cdot u^{-2} + \partial T(z_2) \cdot u^{-1} + {:}TT{:}(z_2) + \cdots$

The Poincar\'e residue of the logarithmic form extracts the coefficient
of $d\log u$ after the chiral product has been expanded in OPE modes.
The channel packet is
\[
T_{(0)}T=\partial T,\qquad
T_{(1)}T=2T,\qquad
T_{(2)}T=0,\qquad
T_{(3)}T=(c/2)\mathbf{1}.
\]
Equivalently, after the logarithmic collision form shifts pole order,
the Virasoro collision kernel has the local normalization
\[
r^{\mathrm{Vir}}(u)=\frac{c/2}{u^3}+\frac{2T}{u}.
\]
The derivative term is the translation/bracket mode.  The scalar
central term is $T_{(3)}T=(c/2)\mathbf{1}$.  Neither term is obtained
by taking an ordinary meromorphic residue of $u^{-4}d\log u$.

The scalar level of the shadow obstruction tower
$\Theta_{\mathrm{Vir}_c}^{\leq 2} = \kappa(\mathrm{Vir}_c) = c/2$
(Theorem~\ref{thm:explicit-theta}) gives the genus-$1$ scalar trace
$F_1(\mathrm{Vir}_c)=\kappa(\mathrm{Vir}_c)/24=c/48$.  The ternary
residue and the genus-$1$ scalar are controlled by the same central
charge, but they live in different gradings.
\end{computation}

\subsection{\texorpdfstring{$m_4$ and higher transfer trees}{m 4 and higher transfer trees}}

\begin{example}[\texorpdfstring{$m_4$}{m_4}: five transfer trees]
\label{ex:m4-two-loop}
The genus-$0$ transferred operation $m_4$ is indexed by the five planar
binary trees with four leaves.  These are the vertices of the
Stasheff pentagon.  A box or two-loop drawing is a physics heuristic;
the algebraic operation is a tree sum unless the modular Feynman
transform has supplied positive-genus vertices.
\end{example}

\begin{computation}[Virasoro \texorpdfstring{$m_4$}{m_4} transfer packet]
\label{comp:v1-virasoro-m4}
\index{Virasoro algebra!$A_\infty$ operations!$m_4$}
For $\mathrm{Vir}_c$, homotopy transfer
(Theorem~\ref{thm:htt}) gives the pentagon packet for
$m_4(T^{\otimes 4})$.

\emph{Setup.}
The homotopy transfer applied to the bar complex
$\B(\mathrm{Vir}_c) \to H^*(\B(\mathrm{Vir}_c))$
produces operations $\tilde{m}_k = \pi \circ T_k \circ \iota^{\otimes k}$,
where $\pi$ is the projection to cohomology, $\iota$ the inclusion,
and $h$ the contracting homotopy.  The operator $h$ becomes a physical
propagator only after a perturbative comparison identifies the
deformation retract with a gauge-fixed BV complex.
For $m_4$, the recursive formula
(Appendix~\ref{app:homotopy-transfer}) gives:
\begin{multline}\label{eq:m4-transfer}
\tilde{m}_4 = \pi \circ \Bigl[
m_2(hm_2(hm_2 \otimes \id)\otimes \id) +
m_2(hm_2(\id \otimes hm_2)\otimes \id) \\
{} + m_2(hm_2 \otimes hm_2) +
m_2(\id \otimes hm_2(hm_2 \otimes \id)) \\
{} + m_2(\id \otimes hm_2(\id \otimes hm_2))
\Bigr] \circ \iota^{\otimes 4}.
\end{multline}

The five terms correspond to the vertices of the Stasheff
associahedron $K_4$ (pentagon), i.e., the $C_3 = 5$ planar binary
trees with $4$~leaves:
\begin{center}
\renewcommand{\arraystretch}{1.3}
\begin{tabular}{cll}
\toprule
$\Gamma$ & Parenthesization & Internal structure \\
\midrule
$\Gamma_1$ & $((T \cdot T) \cdot T) \cdot T$ & left comb \\
$\Gamma_2$ & $(T \cdot (T \cdot T)) \cdot T$ & left-right-left \\
$\Gamma_3$ & $(T \cdot T) \cdot (T \cdot T)$ & balanced \\
$\Gamma_4$ & $T \cdot ((T \cdot T) \cdot T)$ & right-left-right \\
$\Gamma_5$ & $T \cdot (T \cdot (T \cdot T))$ & right comb \\
\bottomrule
\end{tabular}
\end{center}
Each tree $\Gamma_i$ has $2$ internal edges (homotopy insertions~$h$)
and $3$ internal vertices (binary products~$m_2$).

The input $T^{\otimes 4}$ has total conformal weight $4 \times 2 = 8$.
The transfer output has weight determined by the two homotopy
insertions and the OPE modes used at the three binary vertices.  The
possible weight-$6$ normally ordered fields include
\begin{equation}\label{eq:m4-weight-analysis}
m_4(T^{\otimes 4}) \;\in\; \mathbb{C}\cdot\partial^4 T \;\oplus\;
\mathbb{C}\cdot \partial^2({:}TT{:}) \;\oplus\;
\mathbb{C}\cdot {:}\partial T \partial T{:} \;\oplus\;
\mathbb{C}\cdot {:}T\partial^2 T{:} \;\oplus\;
\mathbb{C}\cdot {:}TTT{:},
\end{equation}
a finite space on which the five tree contributions land.  The central
charge enters through the fourth-order OPE mode
$T_{(3)}T=(c/2)\mathbf{1}$.  Its contribution is a curvature-mode
contribution in the bar differential, not a literal product of two
one-loop amplitudes.

The $A_\infty$ relation at degree~$4$ reads
(Appendix~\ref{app:homotopy-transfer},
equation~\eqref{eq:ainfty-relation}):
\begin{equation}\label{eq:pentagon-ainfty}
\sum_{\substack{r+s+t = 4 \\ s \geq 2}} (-1)^{rs+t}\,
\tilde{m}_{r+1+t}\bigl(\id^{\otimes r} \otimes \tilde{m}_s
\otimes \id^{\otimes t}\bigr) = 0.
\end{equation}
The five terms are:
\begin{center}
\renewcommand{\arraystretch}{1.2}
\begin{tabular}{lccc}
\toprule
$(r,s,t)$ & Operation & Sign $(-1)^{rs+t}$ \\
\midrule
$(0,3,1)$ & $\tilde{m}_2(\tilde{m}_3 \otimes \id)$ & $-1$ \\
$(1,3,0)$ & $\tilde{m}_2(\id \otimes \tilde{m}_3)$ & $-1$ \\
$(0,2,2)$ & $\tilde{m}_3(\tilde{m}_2 \otimes \id \otimes \id)$ & $+1$ \\
$(1,2,1)$ & $\tilde{m}_3(\id \otimes \tilde{m}_2 \otimes \id)$ & $-1$ \\
$(2,2,0)$ & $\tilde{m}_3(\id \otimes \id \otimes \tilde{m}_2)$ & $+1$ \\
\bottomrule
\end{tabular}
\end{center}

These correspond to the five \emph{edges} of the pentagon $K_4$
(not the five vertices, which enumerate the trees for~$m_4$).
Substitution of the transferred Virasoro operations gives cancellation
with the displayed signs.  The cancellation is the pentagon identity;
the exact scalar coefficients depend on the chosen transfer
normalization and are not used as independent theorem data here.
\end{computation}

\begin{theorem}[Tree-level \texorpdfstring{$m_k$}{m_k} structure; \ClaimStatusProvedElsewhere]
\label{thm:mk-tree-level}
\index{tree-level!Feynman}
\textup{(}\emph{quadrant} Open, \emph{presentation} transferred chart
operations, \emph{level} 1 of the Beilinson tower,
\emph{hypothesis package} deformation retract $(H, \iota, \pi, h)$ with
$h m_1 + m_1 h = \mathrm{id} - \iota\pi$; planar binary trees;
genus~$0$.\textup{)}
Given a contraction $(H,\iota,\pi,h)$ of an $A_\infty$ algebra
$\mathcal A$, the transferred operation $m_k^H$ is a finite sum over
planar rooted trees with $k$ leaves:
\[
m_k^H =
\sum_{T \in \mathsf{PT}_k}
\pm\,\pi\,\mu_T(h,\mu;\iota,\ldots,\iota),
\]
with the signs determined by the Koszul rule and the planar order.
Every tree has graph loop number~$0$.
\end{theorem}

\begin{proof}
This is the standard homotopy transfer tree formula
(Kadeishvili; Loday--Vallette; Appendix~\ref{app:homotopy-transfer}).
A planar binary tree with $k$ leaves has $k-1$ internal vertices and
$k-2$ internal edges, hence
$\ell=(k-2)-(k-1)+1=0$.
\end{proof}

\begin{theorem}[Formal all-genus stable-graph expansion; \ClaimStatusProvedHere]
\label{thm:mk-general-structure-vol1}
\index{Feynman transform|textbf}
\textup{(}\emph{quadrant} Open, \emph{presentation} bar / twisting
coalgebra in the modular Feynman-transform package,
\emph{level} 2 of the Beilinson tower with descent to clutching at
level~5, \emph{hypothesis package} pronilpotent completion of
Theorem~\ref{thm:stable-graph-pronilpotent-completion}; finite stable
graphs $\mathsf{Gr}^{\mathrm{st}}_{g,k}$ at fixed $(g, k)$;
Euler-characteristic filtration.\textup{)}
In the genus-completed modular bar construction, completed with respect
to the Euler-characteristic filtration, the operation with $k$ external
inputs decomposes formally as the stable-graph sum
\[
m_k^{(g)} =
\sum_{\Gamma \in \mathsf{Gr}^{\mathrm{st}}_{g,k}}
\frac{1}{|\mathrm{Aut}(\Gamma)|}\,\mu_\Gamma,
\]
where
$g=\sum_v g_v+h^1(\Gamma)$, the vertices carry the genus-$g_v$
modular chiral operations, and the internal edges are the clutching or
homotopy-contraction maps of the modular operad.  For fixed $(g,k)$
the graph set is finite; over all genera the sum is an element of the
pronilpotent completion of
Theorem~\textup{\ref{thm:stable-graph-pronilpotent-completion}}.  No
analytic convergence of perturbative integrals and no identification
of edges with Wick kernels is asserted.  In every quotient by the
Euler-characteristic filtration only finitely many graphs remain, so
the equality is a completed algebraic identity, not an all-graphs
numerical summation.  The genus-$0$, $h^1(\Gamma)=0$ specialization is
the tree formula of
Theorem~\textup{\ref{thm:mk-tree-level}}.
Here $m_k^{(g)}$ is a component of the modular Feynman-transform
package, not a genus-$0$ transferred $A_\infty$ operation relabelled by
a physical loop order.
\end{theorem}

\begin{proof}
By Theorem~\ref{thm:prism-higher-genus} (\ClaimStatusProvedHere), the
genus-$g$ bar complex is governed by the Feynman transform of the
commutative modular operad
(Theorem~\ref{thm:bar-modular-operad}).  By definition of the Feynman
transform (Definition~\ref{def:feynman-transform}; Getzler--Kapranov
\cite{GeK98}), its arity-$k$ component is the stable-graph colimit
over graphs with $k$ legs and total genus~$g$.  Vertices contribute the
corresponding modular chiral operations and edges contribute the
operadic clutching/contraction maps.  The finiteness and the passage
to the completed all-genus product are exactly
Theorem~\ref{thm:stable-graph-pronilpotent-completion}.  Passing from
the ordered lift to the symmetric modular quotient divides by the graph
automorphism group.  For $g=0$ and $h^1(\Gamma)=0$ the stable graphs
are precisely trees.
\end{proof}

\begin{remark}[Formal graph sums and scalar projections]
\label{rem:feyn-formal-graph-scalar}
The Feynman graph loop number $\ell = h^1(\Gamma)$ is distinct from the
stable-curve genus $g=\sum_v g_v+\ell$.  Identifying the edge
contractions with analytic Green kernels, or interpreting the stable
graph sum as a perturbative Feynman integral, is additional physics
structure.  That structure is part of Conjecture~\ref{conj:v1-bar-worldline},
not Theorem~\ref{thm:mk-general-structure-vol1}.  After applying the
scalar trace, Theorem~\ref{thm:multi-weight-genus-expansion} gives
\[
F_g(\cA)=F_g^{\mathrm{sc}}(\cA)
+\delta F_g^{\mathrm{cross}}(\cA),
\qquad
F_g^{\mathrm{sc}}(\cA)=\kappa(\cA)\lambda_g^{\mathrm{FP}},
\]
with $\delta F_1^{\mathrm{cross}}=0$ universally and
$\delta F_g^{\mathrm{cross}}=0$ on the scalar lane of
Definition~\ref{def:scalar-lane}.  A stable graph contribution is
therefore not automatically the scalar Hodge-line term
$\kappa(\cA)\lambda_g$; mixed-channel graphs are exactly the
cross-channel correction.
\end{remark}

\section{\texorpdfstring{BPHZ renormalization recursion from $A_\infty$ relations}{BPHZ renormalization recursion from A infty relations}}
\label{sec:bphz-ainfinity}

\subsection{\texorpdfstring{The $A_\infty$ relations as recursion formula}{The A infty relations as recursion formula}}

\begin{conjecture}[BPHZ recursion = \texorpdfstring{$A_\infty$}{A-infinity} consistency; \ClaimStatusConjectured]
\label{conj:v1-bphz-recursion}
\textup{(}\emph{quadrant} Open, \emph{presentation} chart $A_b$
\textup{(}homotopy edges identified with propagators after BV/bar
comparison\textup{)}, \emph{level} 1 of the Beilinson tower,
\emph{hypothesis package} BV/bar comparison morphism identifying
homotopy $h$ with a gauge-fixed propagator kernel; FM boundary faces
identified with subdivergences; uncurved $\Ainf$ stable
vacuum.\textup{)}
The \emph{uncurved} $A_\infty$ relations (with $m_0 = 0$, corresponding to a stable vacuum):
\[\sum_{\substack{r+s+t=n \\ s \geq 1}} (-1)^{rs+t} m_{r+1+t}(\mathrm{id}^{\otimes r} \otimes m_s \otimes
\mathrm{id}^{\otimes t}) = 0\]

match the Bogoliubov--Parasiuk--Hepp--Zimmermann recursion relations
for renormalized Feynman amplitudes after a BV/bar comparison identifies
bar homotopy edges with perturbative propagators.  In the curved case
$m_0\neq 0$, the additional $s=0$ terms correspond to tadpole or vacuum
insertions; they are absent in the uncurved stable-vacuum BPHZ surface.

In that comparison, the $n$-th order amplitude has the recursive shape
\[\mathcal{A}^{(n)} = -\sum_{\substack{\text{proper subgraphs} \\ \Gamma' \subset \Gamma}}
\frac{1}{|\text{Aut}(\Gamma')|} \cdot \mathcal{A}^{(<n)}(\Gamma') \cdot
\mathcal{A}_{\text{reduced}}(\Gamma/\Gamma')\]

where the sum is over all ways to factor $\Gamma$ into a lower-order subgraph $\Gamma'$
and the reduced graph $\Gamma/\Gamma'$.
\end{conjecture}

\begin{remark}[Cubic boundary identity]
At arity~$3$, the $A_\infty$ relation is
\begin{align*}
&m_1(m_3(\phi_1 \otimes \phi_2 \otimes \phi_3))
+ m_2(m_2(\phi_1 \otimes \phi_2) \otimes \phi_3)
+ m_2(\phi_1 \otimes m_2(\phi_2 \otimes \phi_3)) \\
&+ m_3(m_1(\phi_1) \otimes \phi_2 \otimes \phi_3) + \cdots = 0.
\end{align*}
The composites $m_i(\cdots m_j \cdots)$ are algebraic
boundary-factorization terms. Under a BV/bar comparison that sends
homotopy edges to propagators and FM boundary faces to graph
subdivergences, they become factorizable subgraph contributions. The
relation says that the total boundary contribution is $m_1$-exact,
which is the algebraic form of the BPHZ consistency condition
$\mathcal{A}_{\text{ren}}(\Gamma)
= \mathcal{A}_{\text{bare}}(\Gamma)
- \sum_{\text{subdivergences}} \mathcal{A}_{\text{counter}}(\Gamma')$.
The signs $(-1)^{rs+t}$ are the Keller/Koszul signs.  In the FM model
they agree with boundary-orientation signs; in a perturbative field
theory they may also combine with fermionic signs, but that is
model-dependent.
\end{remark}

\begin{example}[One-loop BPHZ shape]
\label{ex:bphz-one-loop}
For a scalar one-loop diagram $\Gamma$ with 3 external legs:

\emph{Bare amplitude:}
\[\mathcal{A}_{\text{bare}}(\Gamma) = \int d^4k \,
\frac{1}{k^2(k-p_1)^2(k-p_1-p_2)^2}\]

This diverges as $k \to \infty$ (UV divergence).

\emph{BPHZ subtraction:}
\[\mathcal{A}_{\text{ren}}(\Gamma) = \mathcal{A}_{\text{bare}}(\Gamma) -
\mathcal{A}_{\text{tree}}|_{\text{evaluated at loop momentum}}\]

The tree-level contribution is:
\[\mathcal{A}_{\text{tree}} = m_2(m_2(\phi_1 \otimes \phi_2) \otimes \phi_3)\]

The $A_\infty$ relation:
\[m_1(m_3(\phi_1 \otimes \phi_2 \otimes \phi_3)) + m_2(m_2(\phi_1 \otimes \phi_2)
\otimes \phi_3) + \cdots = 0\]

tells us:
\[m_3(\phi_1 \otimes \phi_2 \otimes \phi_3) = -h\bigl(m_2(m_2(\cdots))\bigr) + \cdots\]
where $h$ is the homotopy from a deformation retract
$h m_1 + m_1 h = \mathrm{id} - \iota\pi$; in a gauge-fixed BV model it
may be represented by a propagator.

The BPHZ recursion has the same formal shape: the renormalized one-loop
amplitude equals the bare amplitude minus the tree-level counterterm.
This paragraph is a shape comparison.  The equality with a bar formula
requires a BV/bar comparison identifying the homotopy $h$ with the
relevant gauge-fixed propagator kernel and the subdivergence faces with
FM boundary strata.
\end{example}

\subsection{Worldline formalism: configuration spaces as Feynman graphs}

\begin{conjecture}[Bar complex and worldline integrals; \ClaimStatusConjectured]
\label{conj:v1-bar-worldline}
\textup{(}\emph{quadrant} Open, \emph{presentation} bar coalgebra
\textup{(}computational level-2 model\textup{)},
\emph{level} 2 of the Beilinson tower,
\emph{hypothesis package} graph $\Gamma$ with labelled vertices;
configuration-space chain on $\overline{C}_k(X)$; matching of bar form
degree with chosen relative cycle; logarithmic-derivative identity
$dG/G = -\eta_{ij}$ for the Cauchy kernel.\textup{)}
For a graph $\Gamma$ with labelled vertices in
$\{1,\ldots,k\}$, the graph-decorated bar form
\[
\omega_\Gamma =
\phi_1(z_1) \otimes \cdots \otimes \phi_k(z_k)
\otimes \bigwedge_{e=(i,j)\in E(\Gamma)}\eta_{ij}
\]
is analogous to the logarithmic part of a worldline Feynman integrand
before integration, when the form degree matches the chosen
configuration-space chain.

The logarithmic forms $\eta_{ij} = d\log(z_i - z_j)$ record the
logarithmic derivatives of Cauchy singularities:
\[
\eta_{ij} = d\log(z_i-z_j).
\]
If $G(z_i,z_j)=(z_i-z_j)^{-1}$, then $dG/G=-\eta_{ij}$; the sign is
part of the chosen propagator normalization.
\end{conjecture}

\begin{evidence}
Compare the graph-dependent bar pairing, defined when the degree and
relative cycle have been specified,
\[
I^{\mathrm{bar}}_\Gamma(\phi_1,\ldots,\phi_k)
= \left\langle
\phi_1(z_1)\cdots\phi_k(z_k)
\bigwedge_{e=(i,j)\in E(\Gamma)}\eta_{ij},
[\overline{C}_k(X),\partial]
\right\rangle,
\]

with the worldline amplitude:
\[
\mathcal{A}_\Gamma(\phi_1,\ldots,\phi_k)
= \int_{C_k(X)} \phi_1(z_1) \cdots \phi_k(z_k)
\cdot \prod_{e=(i,j)\in E(\Gamma)} G(z_i,z_j)\,
dz_1 \cdots dz_k .
\]

The connection is logarithmic rather than literal:
\[
dG/G=-d\log(z_i-z_j)
\]
for the Cauchy kernel.  The FM compactification resolves UV collision
faces; any IR regularization belongs to the global curve or moduli
problem, not to the local collision blow-up alone.
\end{evidence}

\section{Three algebraic constraints behind the Feynman picture}
\label{sec:feynman-anchors}

\begin{remark}[Three algebraic constraints behind the Feynman picture]
\label{rem:feynman-diagrams-anchors}
\index{structural anchors!Feynman diagrams}
\index{shadow-Feynman dictionary!loop quadrichotomy}
The Feynman-diagram reading is governed by three algebraic constraints.
\begin{enumerate}[label=\textup{(\alph*)}]
\item \emph{Universal conductor.} The genus-$1$ free energy
 $F_1=\kappa(\cA)/24$ is the scalar projection of
 $\mathrm{obs}_1(\cA)=\kappa(\cA)\lambda_1$.  In a BV model it is
 seen as a one-loop anomaly term; algebraically it is the genus-$1$
 Hodge-line contribution of Chapter~\ref{chap:universal-conductor}.
\item \emph{Climax identity.} The bar differential is the signed sum
 of collision residues of the chiral product, with
 $\eta_{ij}=d\log(z_i-z_j)$ supplying the Arnold logarithmic form:
 $d_{\mathrm{bar}} = \mathrm{KZ}^*(\nabla_{\mathrm{Arnold}})$
 of Chapter~\ref{ch:climax-theorem}.  This is an algebraic residue
 theorem.  The off-shell/on-shell reading is the separate conjectural
 pairing of Conjecture~\ref{conj:physical-pairing}.
\item \emph{Shadow-Feynman dictionary.} The
 $A_\infty$ operations $m_k$ are stable-tree and stable-graph
 operations.  The shadow coefficients $S_r$ are degree-$r$
 stable-vertex coefficients in the MC tower
 (Remark~\ref{rem:feynman-shadow-loop-dictionary} of
 Chapter~\ref{ch:v1-feynman}); a physical loop interpretation is
 model-dependent.  The shadow-class
 quadrichotomy $\mathsf{G}/\mathsf{L}/\mathsf{C}/\mathsf{M}$
 controls algebraic depth: classes~$\mathsf{G}$,
 $\mathsf{L}$, $\mathsf{C}$ have finite shadow towers, while
 class~$\mathsf{M}$ has infinite algebraic depth with exponential
 base~$C_\cA$
 (Virasoro: $C_{\mathrm{Vir}} = 6$; principal $W_3$ on the
 W-line: $C_{W_3}^{W\text{-line}} = 12$).
\end{enumerate}
The Feynman picture is therefore a controlled visualization of the
stable-graph MC tower, not a replacement for the bar differential or
the shadow-depth definitions.
\end{remark}

\label{sec:summary-feynman-unity}

\section{Feynman / BV diagrams and the K3 1-loop output}
\label{sec:feyn-supplement}

\begin{remark}[Paramodular $1$-loop scalar target at $K3 \times T^2$; \ClaimStatusHeuristic]
\label{rem:feyn-paramodular-1loop}
The Vol~I Feynman-diagram framework specialises at $K3 \times T^2$ to the twisted
11D supergravity $1$-loop diagram: $24$ M5-brane-loop vertices at the Kodaira $I_1$
fibres of the elliptic K3, each contributing a local $1$-loop determinant. These data
define the Igusa scalar target. If the determinant descends to a scalar
automorphic factor on $\mathcal A_2$, and if the local M5 contributions, bulk
anomaly class, character, divisor, and Weyl-vector normalizations match the
Gritsenko--Nikulin Borcherds product, the trivialising modular form is
$\Delta_5$ of weight~$5$. Uniqueness of the Deligne-cohomology trivialiser in
$H^2(\overline{\mathcal A_2},\mathcal O^*)$ requires the corresponding
line-bundle and divisor comparison; the heuristic claim here is the scalar anomaly
target. Four duality frames (Harvey--Moore $1996$; DVV
$1996$--$1997$; DMVV $1997$; Witten $1995$) point to the same genus-$2$ automorphic
form, conditional on these comparison data.
\end{remark}

\begin{remark}[Bi-based factorisation for Feynman diagrams on $K3 \times E$;
\ClaimStatusHeuristic]
\label{rem:feyn-bi-based}
The proposed Feynman-diagram factorisation datum at $K3 \times E$ is bi-based: chiral base
$\mathrm{Ran}(E^{\mathrm{nod,sm}}_{24})$ with nearby-cycles $\mathcal{D}$-modules at the
$24$ nodes; parameter base $\overline{\mathcal{A}_2}$ with Francis--Gaitsgory
$\infty$-factorisation on holonomic $\mathcal{D}$-modules regular-singular on Humbert
$H_1 \cup H_4$; and a candidate period-map composite
$\mathrm{av} = \mathrm{Torelli}_{K3} \circ
\mathrm{KugaSatake} \circ j_{\mathrm{Kodaira}}$.  This last display is
a target comparison of period maps, not a proved equality of
factorisation functors.
\end{remark}

\section{Refined \texorpdfstring{$\Omega$}{Omega} and the explicit 1-loop diagram}
\label{sec:feyn-supplement-ii}

\begin{remark}[Refined 1-loop determinant in refined $\Omega$-background; \ClaimStatusConjectured]
\label{rem:feyn-W15-refined-1loop}
Place twisted 11D supergravity on $\mathbb R^3 \times \mathrm{K3} \times
\mathbb C^2$ with refined $\Omega$-background $(\epsilon_1, \epsilon_2)$
acting by phase rotations on the two $\mathbb C$-planes of the $\mathbb
C^2$ fibre. The formal 1-loop determinant ansatz decomposes into a
Kaluza--Klein (KK) mode sum along the $\mathrm{K3} \times \mathbb C^2$
direction:
\begin{equation}
\label{eq:feyn-W15-1loop-logZ}
\log Z^{(1)}(\tau_{\mathrm K3}, \epsilon_1, \epsilon_2)
 \;=\; -\!\!\sum_{n \in \mathrm{KK}}\!\!
 \log \det\bigl(\Box_n + \epsilon_1\epsilon_2\bigr),
\end{equation}
where the KK mode index runs over
$\mathrm{KK} \;\cong\; H^*(\mathrm{K3}, \mathbb Z) \otimes_{\mathbb Z}
\Lambda^\bullet_{\mathbb C^2}$,
the tensor of the integral K3 cohomology lattice (signature $(3, 19)$,
rank $24$) with the exterior algebra on the complexified $\mathbb C^2$
fibre. The product $\epsilon_1\epsilon_2$ is the refined $\Omega$ mass
term; when a square root is chosen we write it as~$\hbar^2$.  The
self-dual equation $\epsilon_1+\epsilon_2=0$ fixes only the ratio
$\epsilon_1/\epsilon_2=-1$; the numerical value
$\hbar^2=-1/8$ is the additional normalization
$\epsilon_1=-\epsilon_2=1/(2\sqrt2)$ used below.

The scalar normalization used for the self-dual target assigns the
$24$ Kodaira $I_1$ fibres a total local contribution
$24\cdot(1/8)=3$ in the Harvey--Moore
$j$-regularised scheme.  Adding the separated fibre comparator
$\chi(\mathcal O_{\mathrm K3})=2$ gives the weight-$5$ target for the
Gritsenko--Nikulin form $\Delta_5$.  This is a normalization and
comparison statement, not a proof that the determinant descends to an
automorphic line.  A Deligne-cohomology anomaly-cancellation statement
in $H^2(\overline{\mathcal A}_2, \mathcal O^*)$ requires the scalar
normalization, anomaly-line descent, character, divisor, and Weyl-vector
comparisons of Remark~\ref{rem:feyn-paramodular-1loop}. Primary: Costello
2021 (arXiv:1610.04144) Sec.~3; Costello--Gaiotto 2018
(arXiv:1812.09257); Nekrasov--Okounkov 2003 (hep-th/0306238);
Harvey--Moore 1996 (Nucl.~Phys.~B463); Gritsenko--Nikulin 1998 Thm~2.1.
\end{remark}

\begin{remark}[Explicit diagrammatic computation; \ClaimStatusConjectured]
\label{rem:feyn-W15-explicit-diagram}
Expanding equation~\eqref{eq:feyn-W15-1loop-logZ} diagrammatically:
each KK mode $n \in \mathrm{KK}$ contributes the one-loop diagram of a
massive scalar with mass$^2 = \epsilon_1\epsilon_2$ on
$\mathbb R^3$ (the topological factor). Integrating out the $\mathbb R^3$
factor via zeta-function regularisation yields
\[
\log\det(\Box_n + \epsilon_1\epsilon_2)
 \;=\; -\zeta'(0; \Box_n + \epsilon_1\epsilon_2)
\]
and the conjectural scalar comparison reorganizes the KK-mode sum as
the Jacobi-form generating function of $\phi_{0,1}^{\mathrm K3}$
coefficients in the BKM denominator product:
\[
\log Z^{(1)}
 \;=\; \sum_{\alpha \in L^{++}}
 c_{\mathrm K3}(\alpha\cdot\alpha)\;
 \log\bigl(1 - e^{-\alpha \cdot v}\bigr)
\]
where $L = H^*(\mathrm{K3}, \mathbb Z)$ is the Mukai lattice and
$L^{++}$ the BKM positive cone. The expected Borcherds-product target is
$\Delta_5(v) = \prod_{\alpha \in L^{++}}
(1 - e^{-\alpha\cdot v})^{c_{\mathrm K3}(\alpha\cdot\alpha)}$ at
self-dual slice. Bruinier Prop.~5.1 (Bruinier 2002, LNM~1780) identifies
the exponents $c_{\mathrm K3}(\alpha\cdot\alpha)$ with Heegner-divisor
multiplicities of the Borcherds lift of $\phi_{0,1}^{\mathrm K3}$.
The Vol.~III K3 one-loop compute suite checks the leading Heegner
labels against EOT 2011 Table~1.
\end{remark}

\begin{remark}[Four-loop diagrammatic expansion on $K3\times\bC^2$;
\ClaimStatusConditional]
\label{rem:feyn-4loop-explicit-expansion}
\index{4-loop Feynman!explicit expansion}
\index{K3 4-loop enumeration}
Under the normalization and BV-to-Heegner comparison of
\texttt{chapters/connections/bv\_brst.tex}, the four-loop BV obstruction
of twisted $11$D~SUGRA on $\bR^3\times K3\times\bC^2$ is represented by
the eleven connected four-loop vacuum topologies tabulated in
Table~\ref{tab:feyn-4loop-topologies}.  The table is a transported
consistency packet from Theorem~\ref*{thm:bvbrst-4loop-feynman-sum};
this chapter does not independently enumerate the four-loop graph
complex.  The symmetry factors $|\mathrm{Aut}(\Gamma)|$ use Kreimer's
four-loop Hopf algebra entries~\cite[\S 3, Table 2]{Kreimer1998} and
the Bogner--Weinzierl tensor-reduction
database~\cite[Appendix B]{BognerWeinzierl2017}.
\begin{table}[ht]
\centering
\small
\begin{tabular}{clrc}
\toprule
$\#$ & Topology $\Gamma$ & $\mathrm{Amp}(\Gamma)/[H_4]$ & pair \\
\midrule
$1$ & Wheel $W_4$ & $+576$ & (a) \\
$2$ & Double theta $\Theta\Theta$ & $+72$ & (b) \\
$3$ & Ladder $\mathrm{Lad}_4$ & $-576$ & (a) \\
$4$ & Box with diagonals $\Box^{\mathrm{diag}}_4$ & $-72$ & (b) \\
$5$ & Two-triangle $\Delta\Delta$ & $-48$ & (c) \\
$6$ & Dumbbell-tadpole $\Theta\!+\!\mathrm{tad}$ & $0$ &; \\
$7$ & Pretzel $\mathcal{P}$ & $+48$ & (c) \\
$8$ & $Z$-snake $Z_4$ & $0$ &; \\
$9$ & Sunrise-square $\mathrm{Sun}_4$ & $+12$ & (d) \\
$10$ & Figure-eight-loop $8\!+\!\mathrm{loop}$ & $0$ &; \\
$11$ & Tadpole-tadpole $T^2$ & $0$ &; \\
\midrule
& total & $+12$ & \\
\bottomrule
\end{tabular}
\caption{The eleven connected four-loop Feynman topologies on
$K3\times\bC^2$ and their amplitudes as multiples of the Humbert class
$[H_4]$. Cancellation pairs: (a) wheel--ladder by CGP open-closed
reflection; (b) $\Theta\Theta$--$\Box^{\mathrm{diag}}$ by Humbert
$H_1\leftrightarrow H_4$ parity; (c) two-triangle--pretzel by
Borcherds singular-theta parity; (d) sunrise-square is the residual
Gritsenko--Nikulin witness of $c_{\phi_{-2,1}}(4) = 12$ at Humbert
discriminant $-4$. Proof:
Vol~I~\cite[Theorem~\ref*{thm:bvbrst-4loop-feynman-sum}]{Lorgat1}
at \texttt{chapters/connections/bv\_brst.tex}
\S\ref*{subsec:bvbrst-4loop-enumeration}.}
\label{tab:feyn-4loop-topologies}
\end{table}
The net sum is $c_4 = +12\cdot[H_4] = c_{\phi_{-2,1}}(4)\cdot[H_4]$,
matching the Gritsenko--Nikulin denominator identity
\cite[Thm.~2.1]{GritsenkoNikulin1998} and the four independent
cross-verification paths (DT, heterotic, DMVV, direct Jacobi-form
expansion). The cross-volume triangle of identifications: Vol~I
(BV-cohomology side), Vol~II (Wilson-line counterterm in the HT
bulk-boundary-line triangle), Vol~III (Cartan-matrix root vector of
$\mathfrak{g}_{\Delta_5}$ at discriminant $-4$).
\end{remark}

\begin{remark}[Diagrammatic moves and the 5-loop extension]
\label{rem:feyn-5loop-companion}
\index{4-loop!diagrammatic moves}
\index{5-loop!cancellation}
The four 4-loop cancellation pairs of
Table~\ref{tab:feyn-4loop-topologies} are exhibited as explicit
diagrammatic moves with sign-verification in Vol~I
\texttt{chapters/connections/bv\_brst.tex}
\S\ref*{subsec:bvbrst-4loop-moves}: pair (a) wheel-ladder by
CGP open-closed reflection $\sigma^{\mathrm{CGP}}$
(Proposition~\ref*{prop:bvbrst-4loop-pair-a}), pair (b)
$\Theta\Theta$-$\Box^{\mathrm{diag}}$ by Humbert
$\iota_{H_1, H_4}$-involution with eigenvalue $-1$ from
Bruinier~\cite[Thm.~2.1]{Bruinier2002}
(Proposition~\ref*{prop:bvbrst-4loop-pair-b}), pair (c)
two-triangle-pretzel by Borcherds singular-theta parity
$\varepsilon_{\mathrm{BST}}$~\cite[\S 10]{Borcherds1998}
(Proposition~\ref*{prop:bvbrst-4loop-pair-c}), and tadpole
vanishing via Costello $\chi(K3)/24$-subtraction
(Proposition~\ref*{prop:bvbrst-4loop-tadpole-vanish}).

The 5-loop extension is a companion calculation, not a new theorem in
this chapter.  Kreimer's Hopf algebra gives~\cite[\S 3,
eq.~(16)]{Kreimer1998} $d_5 = 33$ connected topologies in the relevant
vacuum class.  The detailed cancellation packets are those of Vol~I
Theorem~\ref*{thm:bvbrst-5loop-sum}.  After the BV-to-Heegner
normalization is imposed, the graph-weighted sum satisfies
\[
c_5 = c_{\phi_{-2,1}}(5)\cdot[H_5] = 0,
\]
because an index-$1$ weak Jacobi form has nonzero discriminant
coefficients only in classes $0$ or $3$ modulo~$4$, while
$5\equiv 1\pmod 4$.
\end{remark}

\begin{remark}[Refinement off self-dual; \ClaimStatusConjectured]
\label{rem:feyn-W15-refinement}
Off the self-dual locus ($\epsilon_1 + \epsilon_2 \neq 0$), the
1-loop determinant refines to
\[
Z^{(1), \mathrm{ref}}(\tau_{\mathrm K3}, \epsilon_1, \epsilon_2)
 \;=\; \prod_{\alpha \in L^{++}}
 \bigl(1 - e^{-\alpha\cdot v} q_{\epsilon_1}^{\ell_1(\alpha)}
 q_{\epsilon_2}^{\ell_2(\alpha)}\bigr)^{c^{\mathrm{ref}}_{\mathrm K3}(\alpha)}
\]
with multiplicities $c^{\mathrm{ref}}_{\mathrm K3}(\alpha)$ drawn from
the refined K3 elliptic genus
$\chi^{\mathrm{ref}}_{\mathrm K3}(\tau, z_1, z_2)$ of weight $0$ and
bi-index $(1, 1)$ in two elliptic variables (Iqbal--Kozcaz--Vafa
2007 Sec.~4). The self-dual slice $\epsilon_1 + \epsilon_2 \to 0$
collapses $z_1 \to z, z_2 \to z$ and recovers
$\chi_{\mathrm K3}(\tau, z) = \phi_{0,1}^{\mathrm K3}$. The
Nekrasov--Shatashvili degeneration of the refined partition function
has leading free energy equal to a Yang--Yang functional for the
$q$-Bethe integrable system; this is a limit of the two-parameter
theory, not the generic refined determinant itself.  Here
$q := \exp(-2\pi i \epsilon_2 / \epsilon_1)$. Primary: Nekrasov--Pestun--
Shatashvili 2013 (arXiv:1312.6689) Sec.~3; Nekrasov--Shatashvili 2009
(arXiv:0908.4052); Iqbal--Kozcaz--Vafa 2007.
\end{remark}

\begin{remark}[Cross-volume comparison for refined diagrams]
\label{rem:feyn-W15-cross-volume}
The refined 1-loop diagram has the following cross-volume comparison.
\begin{itemize}
\item \textbf{Vol.~I (here):} the refined Feynman datum is the
 paramodular anomaly-cancellation condition on
 $H^2(\overline{\mathcal A}_2, \mathcal O^*)$. The bar-side collision
 forms are still the $\eta_{ij} = d\log(z_i - z_j)$ one-forms.  The
 perturbative propagators in the determinant are the inverse kinetic
 operators $(\Box_n+\epsilon_1\epsilon_2)^{-1}$; vertices receive a
 refined $(\epsilon_1, \epsilon_2)$-twist via the $\mathbb C^\times
 \times \mathbb C^\times$ rotation of the $\mathbb C^2$ fibre.
\item \textbf{Vol.~II:} the level-3 derived chiral centre
 $Z^{\mathrm{der}}_{\mathrm{ch}}(\cA)$ acts on the level-1
 boundary $\cA$ through the chiral Swiss-cheese pair
 (Theorem~\ref{thm:swiss-cheese-chiral-DT-feynman}), realising the
 algebraic side of the $3$d holomorphic-topological sigma model
 $\bR_{\geq 0} \times C$. The refined $\Omega$ deformation is a
 two-parameter deformation of this $\SCchtop$-action of the closed
 colour on the open colour; the self-dual slice is the unrefined
 topologisation. The dimensional uplift here is the Swiss-cheese
 mechanism, not bar iteration on $\cA$.
\item \textbf{Vol.~III:} the refined CY-$3$ output of the two-stage
 functor $\Phi_3^{(\Sigma_2, C)} = \SpCh_{\Sigma_2, C} \circ
 \PhiFA_3$ at reference curve $C = E_\tau$
 \textup{(}Stage~1 $\PhiFA_3$ canonical $E_3$-holomorphic
 factorisation algebra on $\mathrm{K3} \times E$; Stage~2
 $\SpCh_{\Sigma_2, E_\tau}$ specialisation to the elliptic
 reference curve via factorisation homology over the $2$-cycle
 $\Sigma_2$\textup{)},
 $\Phi_3^{(\Sigma_2, E_\tau)}(D^b\mathrm{Coh}(\mathrm{K3} \times E))$,
 is conjecturally compared with a refined BKM algebra whose denominator
 is the refined Borcherds-lift generating function. The
 $(\epsilon_1, \epsilon_2)$ equivariant data on $\bC^2$ lift via
 Calabi--Yau transport to the normal bundle
 $N_{\mathrm{K3}/\mathrm{K3} \times E}$.
\end{itemize}
The K3 bialgebra is native $E_1$ on $E$ and derived $E_2$ on K3
itself.  The refinement specialises to the Nekrasov $\Omega$-background
on toric $\mathbb C^2$ fibres, not on arbitrary CY-fibred geometries.
The product parameter satisfies
$\hbar^2 = \epsilon_1\epsilon_2$, with
$\hbar_1 = \epsilon_1$ and $\hbar_2 = \epsilon_2$ kept distinct.
Cross-ref Vol.~II
core/frontier companion \texttt{ordered\_associative\_chiral\_kd.tex}
Remarks~\ref{V2-rem:v2-oackd-W15-two-param},
\ref{V2-rem:v2-oackd-W15-e1-e2};
Vol.~III \texttt{k3\_chiral\_bialgebra\_platonic.tex} CY-2 $[2]$-shift.
\end{remark}

\begin{remark}[Refined Feynman summary]
\label{rem:feyn-W15-summary}
The K3 1-loop diagram at the self-dual $\Omega$-locus
$\epsilon_1 + \epsilon_2 = 0$ with $\hbar^2 = -1/8$ has conditional
scalar target $\Delta_5$, with weight target $5 = 2 + 3$ after the
normalizations of Remark~\ref{rem:feyn-paramodular-1loop}. The refinement to generic
$\Omega$ is as follows: the explicit 1-loop Feynman determinant
\eqref{eq:feyn-W15-1loop-logZ} is a KK-mode sum over
$H^*(\mathrm{K3}, \mathbb Z) \otimes \Lambda^\bullet_{\mathbb C^2}$ with
mass term $\epsilon_1\epsilon_2$.  Under the refined BKM recognition
hypotheses this sum exponentiates to the refined Borcherds product in
the refined K3 elliptic genus.  Bruinier
Prop.~5.1 identifies the multiplicities with
$c_{\mathrm K3}(4nm-l^2)$, and the Vol.~III compute suite verifies the
listed self-dual coefficients.  The refinement off self-dual is open at
the rigorous level; the expected control mechanism is the refined
topological vertex.
\end{remark}

\section{Refined \texorpdfstring{$\Omega$}{Omega} loop-by-loop and \texorpdfstring{$\Omega$}{Omega}-intermediary spectral parameters}
\label{sec:feyn-supplement-iii}

\begin{remark}[Loop-by-loop structure of refined $\Omega$-partition function; \ClaimStatusConjectured]
\label{rem:feyn-W16-loop-structure}
The refined $\Omega$-partition function
$Z_{\mathrm{ref}}^{\mathrm{K3}}(\epsilon_1, \epsilon_2)$ on
$\mathbb R^3 \times \mathrm{K3} \times \mathbb C^2$ decomposes
formally in a chosen refined square root $\hbar$ of the mass product
$\epsilon_1\epsilon_2$:
\begin{equation}
\label{eq:feyn-W16-loop-decomposition}
\log Z_{\mathrm{ref}}^{\mathrm{K3}}(\epsilon_1, \epsilon_2)
 \;=\; \log Z^{(0)} + \hbar\, F^{(1)} + \hbar^2\, F^{(2)}
 + \hbar^3\, F^{(3)} + O(\hbar^4),
\end{equation}
This is an auxiliary grading for the refined comparison, not the
standard physical genus power $\hbar^{2g-2}$ and not an analytic
small-$\hbar$ convergence theorem.  With that convention the four
structural terms are:
\begin{itemize}
\item \textbf{Tree ($\hbar^0$).} $Z^{(0)}$ is the classical paramodular
 scalar target. At the self-dual slice $\epsilon_1 + \epsilon_2
 = 0$, the expected target is the Gritsenko--Nikulin Igusa form
 $\Delta_5$ of weight $5$, once the anomaly line, character, divisor,
 and root normalizations have been matched. Primary: Gritsenko
 1999 (alg-geom/9606023); Nikulin 1979 (Izv.~Akad.~Nauk).
\item \textbf{1-loop ($\hbar^1$).} $F^{(1)}$ is the Bruinier-reciprocity
 Heegner-Chern-class contribution. At the normalized self-dual slice
 $\hbar^2=-1/8$, Bruinier 2002 (LNM~1780) Prop.~5.1 identifies the
 Borcherds-lift exponent $c_{\mathrm K3}(4nm-l^2)$ with the
 Heegner-divisor multiplicity of the K3 elliptic genus
 $\phi_{0,1}$.  Bruinier's theorem supplies the multiplicity
 identification; it does not by itself fix the $\Omega$-normalization.
 Primary: Bruinier 2002 Prop.~5.1; Borcherds 1998 JAMS~11.
\item \textbf{2-loop ($\hbar^2$).} On the holomorphic-twist comparison
 locus the target coefficient is $F^{(2)}=0$.  The cited physical input
 is the Sen non-renormalisation theorem for the relevant supersymmetric
 amplitude and Costello's holomorphic-twist cocycle calculation; the
 equality with the refined K3 determinant requires the same comparison
 data as the one-loop target.  The identity $c_2 = 0$ is therefore a
 cancellation target, not an independent axiom. Primary: Sen 2007
 Thm.~3.1; Costello 2021 (arXiv:1610.04144) Sec.~3.
\item \textbf{3-loop ($\hbar^3$).} The target nonzero term carries the
 pentagon-heptagon cocycle obstruction:
 \[
 \phi^{(3)} =
 \zeta(3)c_{\mathrm{symm}} + (25/3)c_{\mathrm{timelike}}
 +(\Phi_{10}^{\mathrm{un}}/\eta^{24})
 c_{\Phi_{10}^{\mathrm{un}}}.
 \]
 This is an off-self-dual recognition target.  It becomes a theorem only
 after the compact Hall source, pairing, parity fixture, completion, and
 associator comparison data of
 Remark~\ref{rem:feyn-W16-self-dual-recovery} have been supplied.  The
 Humbert $H_1$ monodromy order $K^\kappa=2c_+=8$ is the predicted
 torsion order of this cocycle. Primary: Maulik--Okounkov 2019 \S16;
 Drinfeld 1990 (Alg.~Anal.~2, Drinfeld associator pentagon).
\end{itemize}
The refined $\Omega$-partition function contains more than a
power series in the product $\epsilon_1\epsilon_2$: after choosing the
refined square root it carries \emph{odd-power} data
($\hbar^1$ Bruinier, $\hbar^3$ pentagon) that are invisible on the
self-dual parity quotient. The Nekrasov 2003 (hep-th/0206161)
Seiberg--Witten prepotential from instanton counting construction
has an NS-limit analogue in the K3 refined setting via
Nekrasov--Shatashvili 2009 (arXiv:0908.4052), where
$\epsilon_2\to0$ produces the Yang--Yang function of the relevant
$q$-Bethe integrable system.
\end{remark}

\begin{remark}[Self-dual slice and the $\Delta_5$ recognition target; \ClaimStatusConjectured]
\label{rem:feyn-W16-self-dual-recovery}
At the self-dual slice $\epsilon_1 + \epsilon_2 = 0$ with
$\hbar^2 = \epsilon_1 \epsilon_2 = -1/8$, the parity quotient of
\eqref{eq:feyn-W16-loop-decomposition} keeps the unrefined tree and
1-loop scalar targets and forgets the $\mathbb Z/2$-odd refined
cocycle components.  The parity-even remnant of the 3-loop target is
the classical $\zeta(3)c_{\mathrm{symm}}$ associator component; it does
not reconstruct the full
$\Phi_{10}^{\mathrm{un}}/\eta^{24}$-twisted Drinfeld associator. The
full associator comparison is accessible only off the self-dual slice.
The BKM algebra
 $\mathbf H_{\Delta_5}$ is a conditional recognition target, not an output of
 the self-dual collapse. A comparison with a Hall--Drinfeld double requires:
 a compact Hall source for $K3\times E$; a Hall pairing compatible with the
 integration or localization theory; a parity and root fixture matching the
 real roots, imaginary roots, and super signs of the $\Delta_5$ denominator;
 a completion in which the double, the dynamical $R$-matrix, and the
 $\Phi_{10}^{\mathrm{un}}/\eta^{24}$ associator are defined; and comparison
 data for the character, Weyl vector, Humbert divisors, and initial
 Fourier--Jacobi coefficients. With these data fixed, the target is an
 isomorphism from the completed refined Hall double to
 $\mathbf H_{\Delta_5}$. Primary: Iqbal--Kozcaz--Vafa 2007;
Nekrasov--Pestun--Shatashvili 2013 (arXiv:1312.6689) Sec.~3.
\end{remark}

\begin{remark}[Spectral-parameter mechanism via $\Omega$-intermediary; \ClaimStatusHeuristic]
\label{rem:feyn-W16-omega-intermediary}
The spectral parameters of the MO-Yangian on
$\mathrm{Hilb}^n(\mathrm K3)$ or on the chiral side do not arise
directly from the normal bundle $N_{\mathrm{K3}/\mathrm{K3} \times E}$.
The mechanism is:
\begin{enumerate}
\item \textbf{Geometric input.} The embedding $\mathrm{K3} \hookrightarrow
 \mathrm{K3} \times E$ has normal bundle
 $N_{\mathrm{K3}/\mathrm{K3} \times E} \cong p^* T_E$ (pullback of
 the tangent bundle of the elliptic curve). Since $c_1(T_E) = 0$,
 the normal bundle is topologically trivial and carries no intrinsic
 spectral parameter.
\item \textbf{$\Omega$-intermediary.} Place a two-parameter
 $\Omega$-background on $E$ with parameters
 $(\epsilon_1^E, \epsilon_2^E) \in
 \mathrm{Lie}(U(1)_A \times U(1)_B)$ rotating the two circle
 factors of $E = \mathbb R^2 / \Lambda$. The $\Omega$-background
 deforms the commutative product on $H^*_T(E)$ to the non-commutative
 Moyal--Weyl product with parameter $\Theta = \epsilon_1^E
 \wedge \epsilon_2^E$.
\item \textbf{Spectral output.} After choosing an overall scale
 $\hbar_E$, the arguments of the R-matrix $R^{\mathrm{MO}}(u, v)$ are
 $u = 2\pi i \epsilon_1^E/\hbar_E$,
 $v = 2\pi i \epsilon_2^E/\hbar_E$.
 The compact K3 side sees only the sum parameter
 $\hbar_E = \epsilon_1^E + \epsilon_2^E$ (the overall $\Omega$-mass
 integrated over $E$),
 not the individual $\epsilon_i^E$; this encodes the compactness
 constraint that eliminates the second $\mathbb{C}^\times$ scaling
 (see Vol~II \S\ref{V2-sec:fsc-K3-vs-A2}).
\end{enumerate}
The chain $N_{C/Y} \to \Omega \to (u, v)$ is
load-bearing and requires $\Omega$ as an intermediary. Skipping the
$\Omega$ step (attempting to derive spectral parameters directly from
$c_1(N_{C/Y})$) fails identically: $c_1(p^* T_E) = 0$, so the direct
route produces zero, not $(\epsilon_1^E, \epsilon_2^E)$. Primary:
Costello--Yagi 2018 (arXiv:1810.01970) \S2--3; Nekrasov 2003
(hep-th/0206161); Nekrasov--Shatashvili 2009 (arXiv:0908.4052)
for the $\Omega$-on-$E$ formulation.
\end{remark}

\begin{remark}[MO-on-$\mathrm{Hilb}(\mathrm K3)$ distinct from $\mathrm{CoHA}(\mathbb A^2)$; \ClaimStatusConditional]
\label{rem:feyn-W16-MO-vs-SV-A2}
The Maulik--Okounkov stable-envelope Yangian
$Y^{\mathrm{MO}}_{\mathrm K3}$ on $\mathrm{Hilb}^n(\mathrm K3)$
(Maulik--Okounkov 2019 Thm.~13.1) is distinct from the
Schiffmann--Vasserot CoHA $\mathrm{CoHA}(\mathbb{A}^2) = Y^+_\hbar
(\widehat{\widehat{\mathfrak{gl}}}_1)$ (Schiffmann--Vasserot 2013,
Publ.~Math.~IHES~118, Thm.~5.1) in the following proved respects,
with the vertex-closure comparison stated as a target:
\begin{itemize}
\item $\mathrm{CoHA}(\mathbb A^2) = Y^+$ is the \emph{positive half}
 of the Yangian (shuffle algebra of symmetric polynomials with
 $(\epsilon_1, \epsilon_2)$-deformed kernel), not the full Yangian;
 the full Yangian is the Drinfeld double
 $\mathcal{D}_\hbar(\mathrm{CoHA}(\mathbb A^2))
 = Y^+ \oplus Y^- \oplus H$.
\item $Y^{\mathrm{MO}}_{\mathrm K3}$ acts on
 $\bigoplus_n H^*_T(\mathrm{Hilb}^n(\mathrm K3))$ with a single
 equivariant parameter $\hbar$; the K3 ambient is compact, with
 no $\mathbb C^\times \times \mathbb C^\times$ scaling available.
 The $\mathbb A^2$ MO-Yangian has two equivariant parameters
 $(\epsilon_1, \epsilon_2)$ and is precisely the quantum toroidal
 algebra $U_{q_1, q_2}(\widehat{\widehat{\mathfrak{gl}}}_1)$.
\item The Lie shadow of $Y^{\mathrm{MO}}_{\mathrm K3}$ is the
 Mukai-lattice Heisenberg $\mathfrak{h}_{\mathrm{Muk}}$ (Nakajima
 1997); the Lie shadow of $\mathrm{CoHA}(\mathbb A^2)$ is the
 positive part of the rank-1 affine $\widehat{\mathfrak{gl}}_1$.
\item The vertex-operator closure of $\mathrm{CoHA}(\mathbb A^2)$
 gives $\mathcal W_{1+\infty}$ at self-dual level
 (Proch\'azka--Rap\v{c}\'ak 2018).  The parallel K3 target is that
 the vertex closure of $Y^{\mathrm{MO}}_{\mathrm K3}$ on K3 Fock
 maps to the Borcherds--Kac--Moody denominator datum
 $\mathfrak{g}_{\Delta_5}$. It recovers the algebra only after the
 root spaces, parity signs, completion, PBW/no-extra-root condition,
 and denominator comparison of
 Remark~\ref{rem:feyn-W16-self-dual-recovery} have been proved.
\end{itemize}
The K3-side and $\mathbb A^2$-side MO-Yangians are therefore not
related by a specialization of equivariant parameters alone: the K3
ambient is compact and the $\mathbb A^2$ ambient is affine with a
two-torus scaling.  They are distinct members of the MO
stable-envelope Yangian class selected by different ambient geometries.
Primary:
Maulik--Okounkov 2019; Schiffmann--Vasserot 2013; Nakajima 1997
Thm.~0.6; Proch\'azka--Rap\v{c}\'ak 2018 (arXiv:1711.06888).
\end{remark}

\section{Nekrasov partition function of $\mathcal{T}[A_1, \Sigma_{0,24}]$ at self-dual $\Omega$}
\label{sec:feyn-nekrasov-A1-sigma-0-24}

The self-dual $\Omega$-slice $\epsilon_1 = -\epsilon_2 = 1/(2\sqrt 2)$
with $\hbar^2 = \epsilon_1\epsilon_2 = -1/8$ is the distinguished
parameter of the K3 $1$-loop Feynman determinant
\eqref{eq:feyn-W15-1loop-logZ}. The class-$\mathcal S$ realisation
used as a comparison channel is the $A_1$ theory on a $24$-punctured
sphere $\mathcal T[A_1, \Sigma_{0,24}]$; its Nekrasov partition
function, the AGT map to Liouville on $\Sigma_{0,24}$, the protected
Schur index identified by Beem--Rastelli with the character of
$\mathcal X(\mathcal T[A_1,\Sigma_{0,24}])$, its conditional finite
comparison with the self-dual Igusa target
$(\Phi_{10}^{\mathrm{un}})^{-1}$, and the Humbert-pole
residue targets predicted by the Jacobi coefficients $c_n =
c_{\phi_{-2,1}}(n)$ are recorded in this section.

\begin{proposition}[Self-dual AGT block on
\texorpdfstring{$\Sigma_{0,24}$}{Sigma 0,24}; \ClaimStatusProvedElsewhere]
\label{prop:feyn-nekrasov-self-dual}
\textup{(}\emph{quadrant} class-$\mathcal{S}$ \textup{(}cited\textup{)},
\emph{presentation} Nekrasov instanton sum compared with Liouville
correlator, \emph{level} 1 of the Beilinson tower
\textup{(}$2d$ chiral algebra side\textup{)}, \emph{hypothesis package}
$24$ distinct marked points on $\Sigma_{0,24}$; AGT normalisation
$c_{\mathrm{Liou}} = 1 + 6Q^2$, $Q = b + b^{-1}$,
$b^2 = -\epsilon_2/\epsilon_1$; symmetric specialisation
$\alpha_i = Q/2$.\textup{)}
Let $\mathcal T = \mathcal T[A_1, \Sigma_{0,24}]$ be the
class-$\mathcal S$ theory of type $A_1$ on the $24$-punctured
sphere with distinct marked points $z_1,\ldots,z_{24}$, placed on the $\Omega$-background
$\mathbb R^4_{\epsilon_1, \epsilon_2}$ at
$\epsilon_1 = -\epsilon_2 = 1/(2\sqrt 2)$, hence
$\hbar^2 = \epsilon_1\epsilon_2 = -1/8$. Then the Nekrasov
instanton partition function (Nekrasov~2003 hep-th/0206161;
Nekrasov~2004 Adv.\ Theor.\ Math.\ Phys.~$7$) equals, under the
AGT normalization conventions, the Liouville conformal block
\[
 Z^{\mathrm{Nek}}_{\mathcal T}(\underline{\alpha};\epsilon_1,\epsilon_2)
 \Bigl|_{\epsilon_1 = -\epsilon_2 = 1/(2\sqrt 2)}
 \;=\;
 \bigl\langle V_{\alpha_1}(z_1)\cdots V_{\alpha_{24}}(z_{24})\bigr\rangle_{
 \mathrm{Liou},\,c_{\mathrm{Liou}} = 25},
\]
with $V_{\alpha_i}(z_i)$ Liouville primaries and, on the symmetric
external-leg specialization used below, external momenta
$\alpha_i = Q/2$, where
$Q = b + b^{-1}$ with $b^2 = -\epsilon_2/\epsilon_1 = 1$ so
$Q = 2$ and $c_{\mathrm{Liou}} = 1 + 6 Q^2 = 25$.
\end{proposition}

\begin{proof}
Alday--Gaiotto--Tachikawa~2009 (arXiv:0906.3219) Thm.~$1$
identifies the instanton partition function of the
class-$\mathcal S$ theory on $\Sigma_{g,n}$ with the Liouville
CFT correlator on the same punctured surface under the AGT
dictionary $c_{\mathrm{Liou}} = 1 + 6 Q^2$,
$Q = b + b^{-1}$, $b^2 = -\epsilon_2/\epsilon_1$. At the
self-dual slice $\epsilon_1 = -\epsilon_2$ we get $b = 1$,
$Q = 2$, $c_{\mathrm{Liou}} = 25$. Setting
$\alpha_i = Q/2$ at every puncture selects the symmetric
external-leg specialization.  Choosing
$z_i=\zeta_{24}^{\,i-1}$ is a cyclotomic marking of the pointed
sphere; it does not produce a geometric Mathieu $M_{24}$ action on
$\mathbb{CP}^1$.  The $M_{24}$ input belongs to the K3 elliptic-genus
moonshine module and enters only in the later finite Igusa-target
comparison. Primary
inputs: Alday--Gaiotto--Tachikawa~2009 Thm.~$1$ (AGT
$Z^{\mathrm{Nek}} = \langle V\cdots V\rangle_{\mathrm{Liou}}$),
Gaiotto~2009 (arXiv:0904.2715 \S~$4$) (class-$\mathcal S$
construction), Eguchi--Ooguri--Tachikawa~2010 (K3
$M_{24}$-symmetry), Nekrasov~2003 Thm.~$4.1$
(instanton-counting partition function). The later K3 one-loop
determinant is only scalar evidence for the conditional $\Delta_5$
target comparison; it is not used in the AGT identification.
\end{proof}

\begin{remark}[Schur index and Beem--Rastelli target comparison]
\label{rem:feyn-nekrasov-schur-limit}
The Schur index of a $4d$ $\mathcal N = 2$ superconformal
theory (Gadde--Rastelli--Razamat--Yan~2011 arXiv:1104.3850)
is a protected specialization of the superconformal index, equivalently
the $S^3\times S^1$ partition function at the Schur fugacity.  It is
not obtained by sending an $\Omega$-parameter difference to zero while
keeping $\epsilon_1\epsilon_2=-1/8$; that would mix two different
backgrounds.  Beem--Rastelli~2013 (arXiv:1312.5344) Thm.~$1$
identifies this Schur index with the graded character of the
$2d$ chiral algebra $\mathcal X(\mathcal T)$ attached to
$\mathcal T$:
\[
 \mathcal I_{\mathrm{Schur}}^{\mathcal T}(q)
 \;=\; \mathrm{ch}_q\bigl(\mathcal X(\mathcal T)\bigr)
 \;=\; \mathrm{Tr}_{\mathcal X}\, q^{L_0 - c/24}.
\]
For $\mathcal T = \mathcal T[A_1, \Sigma_{0,24}]$ with a chosen
K3-moonshine marking, this theorem supplies the $2d$ chiral algebra
$\mathcal X(\mathcal T)$ attached to the $4d$ theory. It does not by itself
identify $\mathcal X(\mathcal T)$ with $\mathbf H_{\Delta_5}$. The intended
physical comparison is conditional: one must construct the proposed
symplectic-boson or lattice realization, fix the $M_{24}$-invariant operation,
match central charge, levels, root grading, and parity with the
$\Delta_5$ denominator algebra, and then run a finite Igusa-target comparison.
Only after that recognition theorem, or an explicit finite coefficient bound
for it, may one assert the reciprocal Igusa character at the Mukai-doubled
slice:
\[
 \mathrm{ch}_q\bigl(\mathcal X(\mathcal T)\bigr)
 \stackrel{?}{=}\frac{1}{\Phi_{10}^{\mathrm{un}}(Z)}
 \Bigl|_{\text{self-dual Siegel}},
\]
as a target comparison with Theorem~\ref{V3-thm:K3xE-coha-character-igusa}
of Vol~III, not as a consequence of Beem--Rastelli alone. Primary:
Beem--Rastelli~2013 Thm.~$1$;
Gadde--Rastelli--Razamat--Yan~2011; Eguchi--Ooguri--
Tachikawa~2010; Gritsenko--Nikulin~1998 ($\Phi_{10}^{\mathrm{un}}$
denominator identity); Cheng--Duncan--Harvey~2014
(arXiv:1208.4074, umbral moonshine comparison data). Cross-ref Vol~III
toric\_cy3\_coha.tex
Theorem~\ref{V3-thm:K3xE-coha-character-igusa}.
\end{remark}

\begin{remark}[Humbert-pole residue targets
\texorpdfstring{$c_n\cdot \lbrack H_n\rbrack$}{c n H n}]
\label{rem:feyn-nekrasov-humbert-residues}
The finite scalar comparison asks for the Nekrasov partition function
$Z^{\mathrm{Nek}}_{\mathcal T[A_1, \Sigma_{0,24}]}$ at
self-dual to have simple logarithmic poles on the Humbert loci
$H_n\subset \overline{\mathcal A}_2$ of Siegel moduli
(Humbert~1900; van der Geer~1988 Ch.~IX). The finite scalar target
comparison predicts that the residue of $\log Z^{\mathrm{Nek}}$ at
$H_n$ is the corresponding Jacobi-form Fourier coefficient:
\[
 \mathrm{Res}_{H_n}\bigl[\log Z^{\mathrm{Nek}}_{\mathcal T}\bigr]
 \stackrel{\mathrm{target}}{=} c_n\cdot [H_n]
 \stackrel{\mathrm{target}}{=} c_{\phi_{-2,1}}(n)\cdot [H_n],
\]
with $c_{\phi_{-2,1}}$ the Fourier coefficient of the weak
Jacobi form $\phi_{-2,1}(\tau, z)$ in the Eichler--Zagier
normalisation $c_{\phi_{-2,1}}(-1) = -1$. This is not a consequence of
Beem--Rastelli alone and does not identify
\(\mathcal X(\mathcal T)\) with \(\mathbf H_{\Delta_5}\). It is the same
conditional finite target comparison as above: AGT supplies the
Liouville/Nekrasov scalar channel, Beem--Rastelli supplies the \(2d\)
chiral algebra \(\mathcal X(\mathcal T)\), and the passage to
\((\Phi_{10}^{\mathrm{un}})^{-1}\) requires the finite Schur--Igusa
recognition theorem, central-charge/level matching, parity/root fixtures,
and the comparison map to the Borcherds denominator. Under those
hypotheses, the Siegel slice of
\(\log(\Phi_{10}^{\mathrm{un}}/\eta^{24})\) expands with
\(c_{\phi_{-2,1}}(n)\) multipliers via the singular-theta lift
(Bruinier~2002 LNM~$1780$ Prop.~$5.1$). In particular, the target value at
$n = 3$ is:
\[
 \mathrm{Res}_{H_3}\bigl[\log Z^{\mathrm{Nek}}_{\mathcal T}\bigr]
 \stackrel{\mathrm{target}}{=} c_3\cdot [H_3]
 \stackrel{\mathrm{target}}{=} -8\cdot [H_3],
\]
checking the scalar $c_3 = -8$ entry of
Theorem~\ref{thm:bvbrst-heegner-humbert-identity}
(bv\_brst.tex eq.~\eqref{eq:bvbrst-heegner-identity}).
Direct $q, y$-expansion of
$\phi_{-2,1} = -(y-2+y^{-1})\prod_{n\geq 1}
(1-q^n y)^2 (1-q^n/y)^2/(1-q^n)^4$ through $q^8$ yields
$c_{\phi_{-2,1}}(3) = -8$, $c_{\phi_{-2,1}}(4) = +12$,
$c_{\phi_{-2,1}}(7) = -39$, $c_{\phi_{-2,1}}(8) = +56$,
which are the target residue coefficients on $H_3, H_4, H_7, H_8$.
Primary: Bruinier~2002 Prop.~$5.1$; Borcherds~1998 JAMS~$11$
Thm.~$10.1$; Eichler--Zagier~$1985$ \S~$9$;
Gritsenko--Nikulin~$1998$ Thm.~$2.1$; bv\_brst.tex
eq.~\eqref{eq:bvbrst-heegner-identity}.
\end{remark}

\begin{remark}[Nekrasov--Shatashvili degeneration and the
\texorpdfstring{$24$}{24}-point Bethe system; \ClaimStatusHeuristic]
\label{rem:feyn-nekrasov-NS-bethe}
The Nekrasov--Shatashvili limit $\epsilon_2 \to 0$ of the
Omega-deformed $\mathcal T[A_1, \Sigma_{0,24}]$ identifies
the Nekrasov free energy as a Yang--Yang functional for a
quantum integrable system with $24$ marked points
(Nekrasov--Shatashvili~2009 arXiv:0908.4052 Thm.~$1$;
Nekrasov--Rosly--Shatashvili~2011 arXiv:1103.3919). This NS
degeneration is adjacent to, not identical with, the fixed
self-dual point $\epsilon_1=-\epsilon_2=1/(2\sqrt2)$.  For the
$A_1$ Gaiotto class-$\mathcal S$
theory on $\Sigma_{0,24}$, the quantum integrable system is
the Gaudin model of type $A_1$ on $\mathbb{CP}^1$ with $24$
marked points (Feigin--Frenkel--Reshetikhin~1994 Comm.\ Math.\
Phys.~$166$; Frenkel~2007 \emph{Langlands
Correspondence\ldots} Ch.~$6$). The Bethe-ansatz spectrum is
the set of solutions $\{u_a\}_{a\geq 1}$ to the Bethe
equations
\[
 \sum_{i=1}^{24}\frac{m_i}{u_a - z_i}
 \;=\;
 \sum_{b\neq a}\frac{2}{u_a - u_b}, \qquad a\geq 1,
\]
with $m_i$ external masses fixed by $\alpha_i = Q/2 = 1$, so
$m_i = 1/2$ for all $i$.  If the marked points are specialised to
$z_i=\zeta_{24}^{i-1}$, the Bethe equations acquire cyclotomic
symmetry.  A Mathieu-moonshine interpretation of the resulting
spectrum would require an additional comparison map from this
cyclotomic Gaudin system to the Cheng--Duncan--Harvey K3 module; it is
not a consequence of the NS theorem. Primary:
Nekrasov--Shatashvili~$2009$ Thm.~$1$; Feigin--Frenkel--
Reshetikhin~$1994$; Frenkel~$2007$ Ch.~$6$;
Cheng--Duncan--Harvey~$2014$ for the moonshine target. Cross-ref
motivic\_shadow\_tower.tex
\S\ref{sec:nekrasov-omega-motivic}.
\end{remark}

\begin{remark}[Cross-volume comparison for the self-dual
Nekrasov construction]
\label{rem:feyn-nekrasov-cross-volume-comparison}
Proposition~\ref{prop:feyn-nekrasov-self-dual} and the
surrounding remarks have the following cross-volume comparison.
\begin{itemize}
\item \textbf{Vol.~I (here):} the Nekrasov partition function
is the $4d$ $\Omega$-deformed instanton sum at self-dual
$\epsilon_1 = -\epsilon_2 = 1/(2\sqrt 2)$. The Feynman
integrand reading is the K3 $1$-loop determinant
\eqref{eq:feyn-W15-1loop-logZ} specialised to the
$24$-term K3 scalar channel; the Humbert residues are target
coefficients controlled by the bar-cobar Koszul obstruction tower
after the finite Igusa-recognition comparison is supplied.
\item \textbf{Vol.~II:} the Liouville-correlator identification is
the $2d$ AGT assembly at level~$1$ of the Beilinson tower.
Its passage to the level-$3$ derived chiral centre via the
chiral Swiss-cheese pair
\textup{(}Theorem~\ref{thm:swiss-cheese-chiral-DT-feynman}\textup{)}
realises the open-closed sigma-model on
$\bR_{\geq 0} \times C$ with $C = \Sigma_{0,24}$; the
$c_{\mathrm{Liou}} = 25$ conformal block is the boundary
\textup{(}open colour\textup{)} datum on which
$\SCchtop$ acts. The Moonshine input is an external K3 module
comparison, not a geometric automorphism of the $24$-punctured
sphere; the dimensional uplift to $3$d HT here is the Swiss-cheese
mechanism, not iterated bar.
\item \textbf{Vol.~III:} the Schur-index/Beem--Rastelli character
supplies the $2d$ chiral algebra to be compared with
$\mathbf H_{\Delta_5}$. The comparison is conditional on Vol~III's
two-stage CY-$3$ functor
$\Phi_3^{(\Sigma_2, C)} =
\SpCh_{\Sigma_2, E^{\mathrm{nod, sm}}_{24}} \circ \PhiFA_3$
at reference curve $C = E^{\mathrm{nod, sm}}_{24}$,
applied to $D^b(K3 \times E)$; Stage~1 producing the canonical
$E_3$-holomorphic factorisation algebra on $K3 \times E$,
Stage~2 specialising to the Kodaira reference curve
$E^{\mathrm{nod, sm}}_{24}$ via factorisation homology over
the $2$-cycle $\Sigma_2$; plus the finite Igusa-target recognition
data of Remark~\ref{rem:feyn-nekrasov-schur-limit}. The $24$-term
structure lives on the $E^{\mathrm{nod, sm}}_{24}$ factorisation
base; see toric\_cy3\_coha.tex
Theorem~\ref{V3-thm:K3xE-coha-character-igusa}.
\end{itemize}
The $4d$ $\Omega$-background is derived $E_2$ on the chiral sector,
whereas the Beem--Rastelli chiral algebra $\mathcal X$ is native
$E_1$.  The self-dual value $\hbar^2=-1/8$ is a distinguished slice,
not the generic two-parameter $\Omega$-background.  The
Beem--Rastelli character is a trace identity, not an algebra
identification.
Primary: Beem--Rastelli~$2013$ Thm.~$1$;
Alday--Gaiotto--Tachikawa~$2009$ Thm.~$1$;
Nekrasov~$2003$ Thm.~$4.1$.
\end{remark}
